\section{Khoảng cách}

\subsection{Đề 1}
\Opensolutionfile{ans}[ans/ans-11-C8-B4]
\caulc

\begin{ex}%[1H8H5-3]
	Cho hình chóp $S.ABCD$ có đáy $ABCD$ là hình vuông tâm $O$, $SA\perp (ABCD)$. Gọi $I$ là trung điểm của $SC$. Khoảng cách từ $I$ đến mặt phẳng $(ABCD)$ bằng độ dài đoạn thẳng nào?
	\choice
	{$IB$}
	{$IA$}
	{$IC$}
	{\True $IO$}
	\loigiai{
		\immini
		{
			Vì $OI\parallel SA$ nên khoảng cách từ $I$ đến mặt phẳng $(ABCD)$ bằng $IO$.
		}{
		\begin{tikzpicture}[scale=0.7, font=\footnotesize,line join=round, line cap=round, >=stealth]
			\path 
			(0,0) coordinate (A)
			++(-140:2) coordinate (B)
			++(0:3.5) coordinate (C)
			($(A)+(C)-(B)$) coordinate (D)
			($(A)+(0,3)$) coordinate (S)
			($(A)!.5!(C)$) coordinate (O)
			($(S)!.5!(C)$) coordinate (I)
			;
			\foreach \i in{B,C,D}{\draw (S)--(\i);};
			\draw (B)--(C)--(D);
			\draw[dashed] (S)--(A)--(B) (A)--(D) (A)--(C) (I)--(O) (B)--(D);
			\pic[draw,angle eccentricity=1.8,angle radius=2mm]{right angle=S--A--D};
			\pic[draw,angle eccentricity=1.8,angle radius=2mm]{right angle=I--O--C};
			\foreach \i/\g in {A/180,B/-90,C/-90,D/0,S/90,O/-90,I/0}
			\fill[black] (\i) circle(1pt)+(\g:3mm)node[scale=1]{$\i$};
		\end{tikzpicture}
		}
	}
\end{ex}
\begin{ex}%[1H8H5-3]
	Cho hình hộp chữ nhật $ABCD.A'B'C'D'$. Khoảng cách giữa hai mặt phẳng $(ABCD)$ và $(A'B'C'D')$ bằng
	\choice
	{$AC'$}
	{$AB'$}
	{$81{m}$}
	{\True $AA'$}
	\loigiai{
		\immini
		{
			Khoảng cách giữa hai mặt phẳng $(ABCD)$ và $(A'B'C'D')$ bằng $AA'$.
		}{
			\begin{tikzpicture}[scale=0.6, font=\footnotesize,line join=round, line cap=round, >=stealth]
				\path 
				(0,0) coordinate (A)
				++(-130:3) coordinate (B)
				++(0:4) coordinate (C)
				($(A)+(C)-(B)$) coordinate (D)
				($(A)!1/2!(C)$) coordinate (O)
				;
				\foreach \i in {A,B,C,D}{
					\coordinate (\i') at ($(\i)+(0,4)$);
				}
				\draw (A')--(B')--(C')--(D')--cycle;
				\draw (B)--(B') (C)--(C') (D)--(D')  (B)--(C)--(D);
				\draw[dashed,thin](B)--(A)--(A') (A)--(D); 
				\foreach \i/\g in {A'/90,B'/90,C'/90,D'/90,A/-90,B/-90,C/-90,D/-90}
				\fill[black] (\i) circle(1pt)+(\g:5mm)node[scale=1]{$\i$};
			\end{tikzpicture}
		}
	}
\end{ex}
\begin{ex}%[1H8H7-2]
	Cho hình chóp $S.ABC$ có đáy $ABC$ là tam giác đều cạnh $a$. Biết $SA\perp (ABC)$ và $SA=a\sqrt{3}$. Tính thể tích khối chóp $S.ABC$.
	\choice
	{$\dfrac{a}{4}$}
	{$\dfrac{a^3}{2}$}
	{\True $\dfrac{a^3}{4}$}
	{$\dfrac{3a^3}{4}$}
	\loigiai{
		\immini
		{
			$V=\dfrac{1}{3}S\cdot h=\dfrac{1}{3}\cdot \dfrac{a^2\sqrt{3}}{4}\cdot a\sqrt{3}=\dfrac{a^3}{4}$.
		}{
			\begin{tikzpicture}[scale=0.5, font=\footnotesize,line join=round, line cap=round, >=stealth]
				\path
				(0,0) coordinate (A)
				++(-60:3) coordinate (B)
				(5,0) coordinate (C)
				($(A)+(0,4)$) coordinate (S)
				;
				\foreach \i in {A,B,C}{\draw (S)--(\i);}
				\draw (A)--(B)--(C);
				\draw[dashed,thin](A)--(C);
				\pic[draw,thin,angle radius=2mm] {right angle = S--A--C};
				\foreach \i/\g in {S/90,A/180,B/-90,C/0}
				\fill[black] (\i) circle(1pt)+(\g:4mm)node[scale=1]{$\i$};
			\end{tikzpicture}
		}
	}
\end{ex}

\begin{ex}%[1H8H5-3]
	Cho hình chóp $S.ABCD$ có đáy là hình bình hành, cạnh bên $SA$ vuông góc với đáy. Biết khoảng cách từ $A$ đến $(SBD)$ bằng $\dfrac{6a}{7}$. Tính khoảng cách từ $C$ đến mặt phẳng $(SBD)$? 
	\choice
	{$\dfrac{12a}{7}$}
	{$\dfrac{3a}{7}$}
	{$\dfrac{4a}{7}$}
	{\True $\dfrac{6a}{7}$}
	\loigiai{
		\immini
		{
			Vì $O$  là trung điểm $AC$ nên\\ $\mathrm{d}(C,(SBD))=\mathrm{d}(A,(SBD))=\dfrac{6a}{7}$.
		}{
			\begin{tikzpicture}[scale=1, font=\footnotesize,line join=round, line cap=round, >=stealth]
				\path 
				(0,0) coordinate (A)
				++(-140:2) coordinate (B)
				++(0:3.5) coordinate (C)
				($(A)+(C)-(B)$) coordinate (D)
				($(A)+(0,3)$) coordinate (S)
				($(A)!.5!(C)$) coordinate (O)
				;
				\foreach \i in{B,C,D}{\draw (S)--(\i);};
				\draw (B)--(C)--(D);
				\draw[dashed] (S)--(A)--(B) (A)--(D)--(B) (A)--(C) (S)--(O);
				\pic[draw,angle eccentricity=1.8,angle radius=2mm]{right angle=S--A--D};
				\foreach \i/\g in {A/-90,B/-90,C/-90,D/-90,S/90,O/-90}
				\fill[black] (\i) circle(1pt)+(\g:3mm)node[scale=1]{$\i$};
			\end{tikzpicture}
		}
	}
\end{ex}
\begin{ex}%[1H8H5-3]
	Cho hình lập phương $ABCD.A'B'C'D'$ có cạnh bằng $a$ (tham khảo hình vẽ). 
	\begin{center}
		\begin{tikzpicture}[scale=0.6, font=\footnotesize,line join=round, line cap=round, >=stealth]
		\path 
		(0,0) coordinate (A)
		++(-130:3) coordinate (B)
		++(0:4) coordinate (C)
		($(A)+(C)-(B)$) coordinate (D)
		($(A)!1/2!(C)$) coordinate (O)
		;
		\foreach \i in {A,B,C,D}{
			\coordinate (\i') at ($(\i)+(0,4)$);
		}
		\draw (A')--(B')--(C')--(D')--cycle;
		\draw (B)--(B') (C)--(C') (D)--(D')  (B)--(C)--(D);
		\draw[dashed,thin](B)--(A)--(A') (A)--(D);
		\foreach \i/\g in {A'/90,B'/90,C'/90,D'/90,A/-90,B/-90,C/-90,D/-90}
		\fill[black] (\i) circle(1pt)+(\g:5mm)node[scale=1]{$\i$};
	\end{tikzpicture}
	\end{center}
	Khoảng cách giữa đường thẳng $BD$ và mặt phẳng $(A'B'C'D')$ bằng
	\choice
	{\True $a$}
	{$a\sqrt{3}$}
	{$3a$}
	{$a\sqrt{2}$}
	\loigiai{
		Vì $BD{ \parallel }B'D'$ nên $BD{ \parallel }(A'B'C'D')$.\\
		Do đó $\mathrm{d}(BD,(A'B'C'D'))=\mathrm{d}(B,(A'B'C'D'))=BB'=a$.\\
		Vì $(ABCD){ \parallel }(A'B'C'D')$ nên $\mathrm{d}((ABCD),(A'B'C'D'))=\mathrm{d}(A,(A'B'C'D'))=AA'=2a\sqrt{3}$}
\end{ex}
\begin{ex}%[1H8H7-4]
	Cho cái nêm hình lăng trụ đứng như hình vẽ.
	\begin{center}
		\begin{tikzpicture}[scale=0.6, font=\footnotesize,line join=round, line cap=round, >=stealth]
			\path 
			(0,0) coordinate (A)
			++(-150:4) coordinate (B)
			++(0:7) coordinate (C)
			($(A)+(C)-(B)$) coordinate (D)
			($(A)!1/2!(C)$) coordinate (O)
			;
			\foreach \i in {A,B}{
				\coordinate (\i') at ($(\i)+(0,2)$);
			}
			\fill[cyan!30] (B)--(A)--(D)--(C)--cycle (A)--(A')--(D)--cycle;
			\draw (B)--(C)--(D)--(A')--(B')--(B) (B')--(C);
			\draw[dashed,thin](B)--(A)--(A') (A)--(D);
			\draw pic[draw, angle radius=2mm]{right angle=C--B--B'};
			\path (B)--(B') node[left,midway]{$7$ cm};
			\path (B)--(C) node[above,midway]{$24$ cm};
			\path (C)--(D) node[below right,midway]{$22$ cm};
		\end{tikzpicture}
	\end{center}
	Thể tích của cái nêm bằng
	\choice
	{$3696\,cm^3$}
	{\True $1848\,cm^3$}
	{$3669\,cm^3$}
	{$1884\,cm^3$}
	\loigiai{
		Diện tích đáy lăng trụ là $S=\dfrac{1}{2}\cdot 7\cdot 24=84\,cm^2$.\\
		Thể tích của khối lăng trụ là $V=S\cdot h=84\cdot 22=1848\,cm^3$.}
\end{ex}
\begin{ex}%[1H8H5-3]
	Cho lập phương $ABCD.A'B'C'D'$ có cạnh bằng $a$ (tham khảo hình bên). 
	\begin{center}
		\begin{tikzpicture}[scale=0.6, font=\footnotesize,line join=round, line cap=round, >=stealth]
		\path 
		(0,0) coordinate (A)
		++(-130:3) coordinate (B)
		++(0:4) coordinate (C)
		($(A)+(C)-(B)$) coordinate (D)
		($(A)!1/2!(C)$) coordinate (O)
		;
		\foreach \i in {A,B,C,D}{
			\coordinate (\i') at ($(\i)+(0,4)$);
		}
		\draw (A')--(B')--(C')--(D')--cycle;
		\draw (B)--(B') (C)--(C') (D)--(D')  (B)--(C)--(D) (A')--(C');
		\draw[dashed,thin](B)--(A)--(A') (D)--(B) (A)--(D);
		\foreach \i/\g in {A'/90,B'/90,C'/90,D'/90,A/-90,B/-90,C/-90,D/-90}
		\fill[black] (\i) circle(1pt)+(\g:5mm)node[scale=1]{$\i$};
	\end{tikzpicture}
	\end{center}
	Khoảng cách giữa hai đường thẳng $BD$ và $A'C'$ bằng
	\choice
	{\True $a$}
	{$a\sqrt{3}$}
	{$\dfrac{a\sqrt{3}}{2}$}
	{$a\sqrt{2}$}
	\loigiai{
		Vì $BD{ \parallel }B'D'$ nên $BD\parallel (A'B'C'D')$.\\
		Do đó $\mathrm{d}(BD,A'C')=\mathrm{d}(BD,(A'B'C'D'))=\mathrm{d}(B,(A'B'C'D'))=BB'=a$.
	}
\end{ex}
\begin{ex}%[1H8H5-3]
	Cho hình chóp $S.ABCD$ có đáy là hình vuông cạnh $a$. Biết $SA=2a$ và $SA$ vuông góc với mặt phẳng $(ABCD)$. Khoảng cách từ điểm $S$ đến mặt phẳng $(ABCD)$ bằng
	\choice
	{$a$}
	{\True $2a$}
	{$3a$}
	{$a\sqrt{2}$}
	\loigiai{
		Vì $SA\perp (ABCD)$ nên $\mathrm{d}(S,(ABCD))=SA=2a$}
\end{ex}
\begin{ex}%[1H8H7-2]
	Cho hình chóp $S.ABC$ có đáy là tam giác đều cạnh $a$. Biết $SA=2a$ và vuông góc với mặt phẳng $(ABC)$. Thể tích của khối chóp $S.ABC$ bằng
	\choice
	{$a^3$}
	{$a^3\sqrt{3}$}
	{$\dfrac{a^3\sqrt{3}}{2}$}
	{\True $\dfrac{a^3\sqrt{3}}{6}$}
	\loigiai{
		Thể tích của khối chóp $S.ABC$ là $V=\dfrac{1}{3}Sh=\dfrac{1}{3}.\dfrac{a^2\sqrt{3}}{4}.2a=\dfrac{a^3\sqrt{3}}{6}$.
	}
\end{ex}
\begin{ex}%[1H8H5-2]
	Cho hình chóp $S.ABC$ có đáy là tam giác đều cạnh $a$. Biết $SA=a\sqrt{3}$ và $SA$ vuông góc với mặt phẳng $(ABC)$. Khoảng cách từ điểm $S$ đến đường thẳng $BC$ bằng
	\choice
	{$a\sqrt{3}$}
	{$a$}
	{\True $\dfrac{a\sqrt{15}}{2}$}
	{$\dfrac{a\sqrt{15}}{3}$}
	\loigiai{
		\immini
		{
			
		Gọi $M$ là trung điểm của $BC$.\\
		Khi đó $\heva{
			& BC\perp AM \\ 
			& BC\perp SA \\ 
		}\Rightarrow BC\perp (SAM)\Rightarrow BC\perp SM$.\\
		Do đó $\mathrm{d}(S,BC)=SM$.\\
		Xét $\Delta SAM$ vuông tại $A$ có $$SM=\sqrt{SA^2+AM^2}=\sqrt{3a^2+{{\left(\dfrac{a\sqrt{3}}{2}\right)}^2}}=\dfrac{a\sqrt{15}}{2}.$$
		Vậy $\mathrm{d}(S,BC)=\dfrac{a\sqrt{15}}{2}$
		}{
			\begin{tikzpicture}[scale=.8, font=\footnotesize,line join=round, line cap=round, >=stealth]
				\path
				(0,0) coordinate (A)
				++(-60:3) coordinate (B)
				(5,0) coordinate (C)
				($(A)+(0,4)$) coordinate (S)
				($(B)!.5!(C)$) coordinate (M)
				;
				\foreach \i in {A,B,C}{\draw (S)--(\i);}
				\draw (A)--(B)--(C) (S)--(M);
				\draw[dashed,thin] (M)--(A)--(C);
				\pic[draw,thin,angle radius=2mm] {right angle = S--A--C};
				\foreach \i/\g in {S/90,A/180,B/-90,C/0,M/-30}
				\fill[black] (\i) circle(1pt)+(\g:4mm)node[scale=1]{$\i$};
			\end{tikzpicture}
	}
}
\end{ex}
\begin{ex}%[1H8H5-3]
	Cho hình chóp tứ giác đều $S.ABCD$ có cạnh đáy bằng $a$, cạnh bên bằng $2a$. Gọi $O$ là tâm của $ABCD$ và $G$là trọng tâm của tam giác $ACD$. Khoảng cách từ điểm $G$ đến mặt phẳng $(SBC)$ bằng
	\choice
	{$\dfrac{a\sqrt{210}}{30}$}
	{\True $\dfrac{2a\sqrt{210}}{45}$}
	{$\dfrac{a\sqrt{210}}{45}$}
	{$\dfrac{a\sqrt{210}}{15}$}
	\loigiai{
		\begin{center}
			\begin{tikzpicture}[scale=1, font=\footnotesize,line join=round, line cap=round, >=stealth]
				\path 
				(0,0) coordinate (A)
				++(-140:3) coordinate (B)
				++(0:4) coordinate (C)
				($(A)+(C)-(B)$) coordinate (D)
				($(A)!1/2!(C)$) coordinate (O)
				($(O)+(0,3.5)$) coordinate (S)
				($(B)!1/2!(C)$) coordinate (M)
				($(S)!(O)!(M)$) coordinate (H)
				;
				\foreach \i in{B,C,D}{\draw (S)--(\i);};
				\draw (B)--(C)--(D) (S)--(M);
				\draw[dashed] (S)--(A)--(B) (C)--(A)--(D)--(B) (S)--(O)--(M) (O)--(H);
				\pic[draw,angle eccentricity=1.8,angle radius=2mm]{right angle=S--O--D};
				\pic[draw,angle eccentricity=1.8,angle radius=2mm]{right angle=O--H--S};
				\foreach \i/\g in {A/180,B/180,C/-90,D/0,O/-70,S/90,M/-90,H/180}
				\fill[black] (\i) circle(1pt)+(\g:3mm)node[scale=1]{$\i$};
			\end{tikzpicture}
		\end{center}
		Vì $\dfrac{\mathrm{d}(G,(SBC))}{\mathrm{d}(O,(SBC))}=\dfrac{GB}{OB}=\dfrac{4}{3}\Rightarrow \mathrm{d}(G,(SBC))=\dfrac{4}{3}\mathrm{d}(O,(SBC))$.\\
		Gọi $M$ là trung điểm của $BC$.\\
		Khi đó $\heva{
			& BC\perp OM \\ 
			& BC\perp SO \\ 
		}\Rightarrow BC\perp (SOM)$.\\
		Trong mặt phẳng $(SOM)$, kẻ $OH\perp SM$.\\
		Khi đó $OH\subset (SOM)$ nên $BC\perp OH$, suy ra $OH\perp (SBC)$ hay $\mathrm{d}(O,(SBC))=OH$.\\
		Ta có $SO=\sqrt{(2a)^2-{{\left(\dfrac{a\sqrt{2}}{2}\right)}^2}}=\dfrac{a\sqrt{14}}{2}$.\\
		Xét $\Delta SOM$ vuông tại $O$ có $$\dfrac{1}{OH^2}=\dfrac{1}{OM^2}+\dfrac{1}{SO^2}=\dfrac{1}{{{\left(\dfrac{a}{2}\right)}^2}}+\dfrac{1}{{{\left(\dfrac{a\sqrt{14}}{2}\right)}^2}}=\dfrac{30}{7a^2}\Rightarrow OH=\dfrac{a\sqrt{210}}{30}.$$
		Vậy $\mathrm{d}(G,(SBC))=\dfrac{4}{3}\mathrm{d}(O,(SBC))=\dfrac{2a\sqrt{210}}{45}$.}
\end{ex}
\begin{ex}%[1H8H5-3]
	Cho hình hộp chữ nhật $ABCD.A'B'C'D'$ có $AB=1$, $BC=2$, $AA'=2$ (tham khảo hình bên).
	\begin{center}
		\begin{tikzpicture}[scale=0.6, font=\footnotesize,line join=round, line cap=round, >=stealth]
			\path 
			(0,0) coordinate (A)
			++(-130:3) coordinate (B)
			++(0:4) coordinate (C)
			($(A)+(C)-(B)$) coordinate (D)
			($(A)!1/2!(C)$) coordinate (O)
			;
			\foreach \i in {A,B,C,D}{
				\coordinate (\i') at ($(\i)+(0,3)$);
			}
			\draw (A')--(B')--(C')--(D')--cycle;
			\draw (B)--(B') (C)--(C') (D)--(D')  (B)--(C)--(D);
			\draw[dashed,thin](B)--(A)--(A') (A)--(D);
			\foreach \i/\g in {A'/90,B'/90,C'/90,D'/90,A/-90,B/-90,C/-90,D/-90}
			\fill[black] (\i) circle(1pt)+(\g:5mm)node[scale=1]{$\i$};
		\end{tikzpicture}
	\end{center}
	Khoảng cách giữa hai đường thẳng $AD'$ và $DC'$ bằng
	\choice
	{$\sqrt{2}$}
	{$\dfrac{\sqrt{6}}{2}$}
	{$\dfrac{2\sqrt{5}}{5}$.}
	{\True $\dfrac{\sqrt{6}}{3}$}
	\loigiai{
		\begin{center}
			\begin{tikzpicture}[scale=0.6, font=\footnotesize,line join=round, line cap=round, >=stealth]
				\path 
				(0,0) coordinate (A)
				++(-130:3) coordinate (B)
				++(0:6) coordinate (C)
				($(A)+(C)-(B)$) coordinate (D)
				($(A)!1/2!(C)$) coordinate (O)
				($(D')!1/2!(C)$) coordinate (M)
				;
				\foreach \i in {A,B,C,D}{
					\coordinate (\i') at ($(\i)+(0,4.5)$);
				}
				\coordinate (M) at ($(D')!1/2!(C)$);
				\draw (A')--(B')--(C')--(D')--cycle;
				\draw (B)--(B') (C)--(C') (D)--(D')  (B)--(C)--(D) (B)--(C')--(D) (C)--(D');
				\draw[dashed,thin](B)--(A)--(A') (A)--(D) (A)--(D') (M)--(B)--(D);
				\foreach \i/\g in {A'/90,B'/90,C'/90,D'/90,A/180,B/-90,C/-90,D/-90,M/0}
				\fill[black] (\i) circle(1pt)+(\g:5mm)node[scale=1]{$\i$};
			\end{tikzpicture}
		\end{center}
		Vì $AD'{ \parallel }BC'$ nên $AD'{ \parallel }(DBC')$.\\
		Do đó $\mathrm{d}(AD',DC')=\mathrm{d}(AD',(DBC'))=\mathrm{d}(A,(DBC'))=\mathrm{d}(C,(DBC'))=h$.\\
		Xét tứ diện $C.BC'D$ có các cạnh $CD,CB,CC'$ đôi một vuông góc nên ta có
		$$\dfrac{1}{h^2}=\dfrac{1}{CB^2}+\dfrac{1}{CD^2}+\dfrac{1}{C{C'}^2}=\dfrac{1}{2^2}+\dfrac{1}{1^2}+\dfrac{1}{2^2}=\dfrac{3}{2}\Rightarrow h=\dfrac{\sqrt{6}}{3}.$$
		Vậy $\mathrm{d}(AD',DC')=\dfrac{\sqrt{6}}{3}$.
	}
\end{ex}


\Closesolutionfile{ans}
\indapan{6}{ans/ans-11-C8-B4}
\cauds
\Opensolutionfile{ans}[ans/ans-11-C8-B4-DS]

\begin{ex}%[1H8V5-3]
	Cho hình chóp $S.  A B C D$ có đáy là vuông cạnh $a$. Biết $S A$ vuông góc với mặt phẳng đáy và $S A=a \sqrt{3}$. Vẽ đường cao $A H$ của tam giác $S A B$.Vẽ đường cao $A K$ của tam giác $S A D$. Khi đó:
	\choiceTF
	{\True $BC\perp AH$}
	{\True Khoảng cách từ $A$ đến mặt phẳng $(S B C)$ bằng: $\dfrac{a\sqrt{3}}{2}$}
	{Khoảng cách từ $A$ đến mặt phẳng $(S B D)$ bằng: $\dfrac{a\sqrt{2}}{7}$}
	{Khoảng cách từ $C$ đến mặt phẳng $(A H K)$ bằng: $\dfrac{a\sqrt{5}}{5}$}
\loigiai{
	\begin{center}
		\begin{tikzpicture}[scale=1, font=\footnotesize,line join=round, line cap=round, >=stealth]
			\path 
			(0,0) coordinate (A)
			++(-140:2) coordinate (B)
			++(0:4) coordinate (C)
			($(A)+(C)-(B)$) coordinate (D)
			($(A)+(0,4)$) coordinate (S)
			($(A)!.5!(C)$) coordinate (O)
			($(S)!(A)!(O)$) coordinate (E)
			($(S)!(A)!(B)$) coordinate (H)
			($(S)!(A)!(D)$) coordinate (K)
			(intersection of S--O and H--K) coordinate (I)
			($(A)!2!(I)$) coordinate (J)
			(intersection of A--J and S--C) coordinate (F);
			;
			\foreach \i in{B,C,D}{\draw (S)--(\i);};
			\draw (B)--(C)--(D) (H)--(F)--(K);
			\draw[dashed] (S)--(A)--(B) (A)--(D) (A)--(C) (D)--(B) (K)--(A)--(H)--(K) (S)--(O) (E)--(A)--(F);
			\pic[draw,angle eccentricity=1.8,angle radius=2mm]{right angle=S--A--D};
			\pic[draw,angle eccentricity=1.8,angle radius=2mm]{right angle=A--E--O};
			\pic[draw,angle eccentricity=1.8,angle radius=2mm]{right angle=A--O--B};
			\foreach \i/\g in {A/-90,B/-90,C/-90,D/-90,S/90,H/180,K/50,F/60,E/0}
			\fill[black] (\i) circle(1pt)+(\g:3mm)node[scale=1]{$\i$};
		\end{tikzpicture}
	\end{center}
\begin{itemchoice}
	\itemch Ta có $\heva{&B C \perp S A \text{ (do } S A \perp(A B C D)) \\& B C \perp A B} \Rightarrow B C \perp(S A B) \Rightarrow B C \perp A H$.
	\itemch  Ta có $B C \perp A H$
	mà $S B \perp A H$ nên $A H \perp(S B C)$ hay $\mathrm{d}(A,(S B C))=A H$.\\
	Tam giác $S A B$ vuông tại $A$ có đường cao $A H$ nên
	$$\dfrac{1}{A H^2}=\dfrac{1}{A B^2}+\dfrac{1}{S A^2} \Rightarrow A H=\dfrac{A B \cdot S A}{\sqrt{A B^2+S A^2}}=\dfrac{a \cdot a \sqrt{3}}{\sqrt{a^2+3 a^2}}=\dfrac{a \sqrt{3}}{2}.$$
	Vậy $\mathrm{d}(A,(S B C))=A H=\dfrac{a \sqrt{3}}{2}$.
	\itemch Gọi $O$ là tâm hình vuông $A B C D$ thì $A O \perp B D$, ta lại có $S A \perp B D$ nên $B D \perp(S A C)$.\tagEX{*}
	\noindent Kẻ đường cao $A E$ của $\Delta S A O$ thì $A E \perp B D$ (do $(*))$.\\
	Vậy $A E \perp(S B D)$ hay $\mathrm{d}(A,(S B D))=A E$.\\
	Ta có $A C=a \sqrt{2}$ (đường chéo hình vuông), suy ra $O A=\dfrac{A C}{2}=\dfrac{a \sqrt{2}}{2}$.\\
	Tam giác $S A O$ vuông tại $A$ có $A E=\dfrac{S A \cdot A O}{\sqrt{S A^2+A O^2}}=\dfrac{a \sqrt{3} \cdot \dfrac{a \sqrt{2}}{2}}{\sqrt{3 a^2+\dfrac{2 a^2}{4}}}=\dfrac{a \sqrt{21}}{7}$.\\
	Vậy $\mathrm{d}(A,(S B D))={AE}=\dfrac{a \sqrt{21}}{7}$.
	\itemch Ta chứng minh được $A K \perp(S C D)$. Khi đó: $\heva{&S C \perp A H \\ &S C \perp A K} \Rightarrow S C \perp(A H K)$.\\
	Gọi $F=S C \cap(A H K)$ thì $S C \perp A F$. 
	Khi đó: $\mathrm{d}(C,(A H K))=C F$.\\
	Ta có $S C=\sqrt{S A^2+A C^2}=\sqrt{3 a^2+2 a^2}=a \sqrt{5}$.\\
	Tam giác $S A C$ vuông tại $A$ có đường cao $A F$ nên:
	$$CF\cdot CS=AC^2\Rightarrow CF=\dfrac{AC^2}{CS}=\dfrac{2a^2}{a\sqrt{5}}=\dfrac{2a\sqrt{5}}{5}.$$
	Vậy $\mathrm{d}(C,(AHK))=CF=\dfrac{2a\sqrt{5}}{5}$.
\end{itemchoice}
}
\end{ex}
\begin{ex}%[1H8V7-2]
Cho hình chóp $S.  A B C$ có mặt bên $(S A B)$ vuông góc với mặt đáy và tam giác $S A B$ đều cạnh $2 a$. Biết tam giác $A B C$ vuông tại $C$ và cạnh $A C=a \sqrt{3}$. Khi đó:
\choiceTF
{\True $S H \perp(A B C)$}
{\True $\mathrm{d}(S,(ABC))=a\sqrt{3}$}
{$\mathrm{d}(C,(SAB))=\dfrac{a\sqrt{3}}{3}$}
{Thể tích của khối chóp $S.  A B C$ bằng $\dfrac{a^3}{6}$}
\loigiai{
	\begin{center}
		\begin{tikzpicture}[scale=1, font=\footnotesize,line join=round, line cap=round, >=stealth]
			\path
			(0,0) coordinate (B)
			++(-45:2) coordinate (C)
			(5,0) coordinate (A)
			($(A)!1/2!(B)$) coordinate (H)
			($(H)+(0,3.5)$) coordinate (S)
			($(A)!(C)!(B)$) coordinate (K)
			;
			\foreach \i in {A,B,C}{\draw (S)--(\i);}
			\draw (A)--(C)--(B);
			\draw[dashed,thin](A)--(B) (S)--(H) (C)--(K);
			\pic[draw,thin,angle radius=2mm] {right angle = S--H--A};
			\pic[draw,thin,angle radius=2mm] {right angle = B--C--A};
			\pic[draw,thin,angle radius=2mm] {right angle = C--K--B};
			\foreach \i/\g in {S/90,A/0,B/180,C/-90,H/-90,K/90}
			\fill[black] (\i) circle(1pt)+(\g:4mm)node[scale=1]{$\i$};
			\path (S)--(B) node[left,midway]{$2a$};
			\path (S)--(A) node[right,midway]{$2a$};
			\path (C)--(A) node[right,midway]{$a\sqrt{3}$};
			\path (B)--(A) node[above,pos=.7]{$2a$};
		\end{tikzpicture}
	\end{center}
\begin{itemchoice}
	\itemch Gọi $H$ là trung điểm $A B$, mà tam giác $S A B$ đều nên $S H \perp A B$. \\
	Ngoài ra $(S A B) \perp(A B C)$ nên $S H \perp(A B C)$.
	\itemch Ta có $\mathrm{d}(S,(A B C))=S H=\dfrac{2 a \cdot \sqrt{3}}{2}=a \sqrt{3}$ (do tam giác $S A B$ đều cạnh $2 a)$.
	\itemch Kẻ đường cao $C K$ của tam giác $A B C$.\\
	Ta có $\heva{&C K \perp A B \\ &C K \perp S H}\Rightarrow C K \perp(S A B) \Rightarrow \mathrm{d}(C,(S A B))=C K$.\\
	Xét tam giác $A B C$ vuông tại $C$ có\\
	$B C=\sqrt{A B^2-A C^2}=\sqrt{4 a^2-3 a^2}=a;  C K=\dfrac{C A \cdot C B}{A B}=\dfrac{a \sqrt{3} \cdot a}{2 a}=\dfrac{a \sqrt{3}}{2}$.\\
	Vậy $\mathrm{d}(C,(S A B))=C K=\dfrac{a \sqrt{3}}{2}$.\\
	\itemch Diện tích đáy hình chóp là: $S_{\Delta A B C}=\dfrac{1}{2} A C \cdot B C=\dfrac{1}{2} a \sqrt{3} \cdot a=\dfrac{a^2 \sqrt{3}}{2}$.
	Thể tích khối chóp là: $\mathrm{V}_{S \cdot A B C}=\dfrac{1}{3} S H \cdot S_{\Delta A B C}=\dfrac{1}{3} \cdot a \sqrt{3} \cdot \dfrac{a^2 \sqrt{3}}{2}=\dfrac{a^3}{2}$
\end{itemchoice}
}
\end{ex}
\begin{ex}%[1H8HV5-4]
Cho hình hộp chữ nhật $A B C D \cdot A’ B’ C’ D’$ có $A B=a, A D=b, A A’=c$. Khi đó:
\choiceTF
{$A B \perp(A D D’ A’)$}
{Khoảng cách từ điểm $A$ đến đường thẳng $B D’$ bằng: $\dfrac{\sqrt{b^2+c^2}}{\sqrt{a^2+b^2+c^2}}$}
{Gọi $I, J$ theo thứ tự là tâm của các hình chữ nhật $A D D’ A’, B C C’ B’$. Khi đó $I J$ là đường vuông góc chung của hai đường thẳng $A D’$ và $B’ C$}
{Khoảng cách hai đường thẳng $A D’$ và $B’ C$ bằng $2a$}
\loigiai{
	\begin{center}
		\begin{tikzpicture}[scale=1, font=\footnotesize,line join=round, line cap=round, >=stealth]
			\path 
			(0,0) coordinate (D)
			++(-130:2) coordinate (A)
			++(0:4) coordinate (B)
			($(D)+(B)-(A)$) coordinate (C)
			($(A)!1/2!(D')$) coordinate (I)
			($(B)!1/2!(C')$) coordinate (J)
			;
			\foreach \i in {A,B,C,D}{
				\coordinate (\i') at ($(\i)+(0,3.5)$);
			}
			\path 
			($(A)!1/2!(D')$) coordinate (I)
			($(B)!1/2!(C')$) coordinate (J)
			($(B)!(A)!(D')$) coordinate (H)
			;
			\draw (A')--(B')--(C')--(D')--cycle;
			\draw (A')--(A)--(B)--(B') (B)--(C)--(C') (B')--(C) (B)--(C');
			\draw[dashed,thin](B)--(A) (C)--(A)--(D)--(B) (C)--(D)--(D') (A)--(D') (D)--(A') (D')--(B) (I)--(J) (A)--(H);
			\pic[draw,angle eccentricity=1.8,angle radius=2mm]{right angle=A--H--D'};
			\foreach \i/\g in {A'/90,B'/90,C'/90,D'/90,A/-90,B/-90,C/-90,D/-90,I/180,J/0,H/50}
			\fill[black] (\i) circle(1pt)+(\g:5mm)node[scale=1]{$\i$};
		\end{tikzpicture}
	\end{center}
\begin{itemchoice}
	\itemch Kẻ đường cao $A H$ trong tam giác $A B D’$, suy ra $\mathrm{d}(A, B D’)=A H$.\\
	Vì $A B C D \cdot A’ B’ C’$ là hình hộp chữ nhật nên $A B \perp(A D D’ A’)$, 
	suy ra $A B \perp A D’$ hay tam giác $A B D’$ vuông tại $A$.
	\itemch Tam giác $A D D’$ vuông tại $D$ có $A D’=\sqrt{A D^2+D D^{\prime 2}}=\sqrt{b^2+c^2}$.\\
	Tam giác $A B D’$ vuông tại $A$ có đường cao $A H$ nên
	$$\dfrac{1}{A H^2}=\dfrac{1}{A B^2}+\dfrac{1}{A D^{\prime 2}} \Rightarrow A H=\dfrac{A B \cdot A D’}{\sqrt{A B^2+A D^{\prime 2}}}=\dfrac{a \sqrt{b^2+c^2}}{\sqrt{a^2+b^2+c^2}}
	$$
	Vậy $\mathrm{d}(A, B D’)=\dfrac{a \sqrt{b^2+c^2}}{\sqrt{a^2+b^2+c^2}}$.
	\itemch Vì $A B C D \cdot A’ B’ C’ D’$ là hình hộp chữ nhật nên $\heva{&A B / / C’ D’ \\ &A B=C’ D’}$\\
	$\Rightarrow A B C’ D’$ là hình bình hành.\\
	Dễ thấy $I, J$ lần lượt là trung điểm của $A D’$ và $B C’$ suy ra $I J$ là đường trung bình của hình bình hành $A B C’ D’ \Rightarrow I J / / A B$, mà $A B \perp A D’$ nên $I J \perp A D’$ \tagEX{1}
	Ta có $A B \perp(B C C’ B’) \Rightarrow A B \perp B’ C \Rightarrow I J \perp B’ C$. \tagEX{2}
	Mặt khác $I J$ cắt cả hai đường thẳng $A D’, B’ C$. \tagEX{3}
	Từ (1), (2), (3) suy ra $I J$ là đường vuông góc chung của hai đường thẳng $A D’$ và $B’ C$. 
	\itemch Ta có $I J=A B=a$.
\end{itemchoice}
}
\end{ex}
\begin{ex}%[1H8HV7-4]
Cho lăng trụ đứng $A B C \cdot A’ B’ C’$ có $A C=a, B C=2 a, \widehat{A C B}=120^{\circ}$. Gọi $M$ là trung điểm của $B B’$. Khi đó: 
\choiceTF
{\True $\mathrm{d}(C{{C}’},(AB{{B}’}{{A}’}))=\dfrac{a\sqrt{21}}{7}$}
{$\mathrm{d}(C{{C}’},AM)=\dfrac{a\sqrt{21}}{12}$}
{\True $A A’ \perp(A B C), A A’ \perp(A’ B’ C’)$}
{\True Biết khoảng cách giữa hai mặt đáy lăng trụ bằng $2 a$. Khi đó thể tích khối lăng trụ là: $a^3\sqrt{3}$}
\loigiai{
	\begin{center}
		\begin{tikzpicture}[scale=0.8, font=\footnotesize,line join=round, line cap=round, >=stealth]
			\path 
			(0,0) coordinate (A)
			++(-40:2.5) coordinate (C)
			(6,0) coordinate (B)
			($(A)!.25!(B)$) coordinate (H)
			($(B')!.5!(B)$) coordinate (M)
			;
			\foreach \i in{A,B,C}{
				\coordinate (\i') at ($(\i)+(0,-4)$);
				\draw (\i)--(\i');};
			\path ($(B')!.5!(B)$) coordinate (M);
			\draw (A')--(C')--(B') (A)--(C)--(B)--cycle (C)--(H);
			\draw[dashed] (A')--(B') (A)--(M);
			\pic[draw,angle eccentricity=1.8,angle radius=2mm]{right angle=A'--A--C};
			\pic[draw,angle eccentricity=1.8,angle radius=2mm]{right angle=C--H--B};
			\pic[draw,"$120^{\circ}$", angle eccentricity=0.6,angle radius=0.6cm]{angle=B--C--A};
			\foreach \i/\g in {A/90,B/90,C/-30,A'/180,B'/0,C'/-90,H/90,M/0}
			\fill[black] (\i) circle(1pt)+(\g:4mm)node[scale=1]{$\i$};
			\path (A)--(A') node[left,midway]{$2a$};
			\path (A)--(C) node[below,midway]{$a$};
			\path (C)--(B) node[below,midway]{$2a$};
		\end{tikzpicture}
	\end{center}
\begin{itemchoice}
	\itemch Ta có $C C’ / / B B’ \Rightarrow C C’ / /(A B B’ A’)$ nên $\mathrm{d}(C C’,(A B B’ A’))=\mathrm{d}(C,(A B B’ A’))$.\\
	Trong mặt phẳng $(A B C)$, kẻ $C H \perp A B$ tại $H$. \tagEX{1}
	Vì $A B C \cdot A’ B’ C’$ là hình lăng trụ đứng nên $A A’ \perp(A B C) \Rightarrow C H \perp A A’$. \tagEX{2}
	Từ (1) và (2) suy ra $C H \perp(A B B’ A’) \Rightarrow \mathrm{d}(C,(A B B’ A’))=C H$. \\
	Xét tam giác $A B C$, có $A B^2=C A^2+C B^2-2 C A \cdot C B \cdot \cos 120^{\circ}=7 a^2$ $\Rightarrow A B=a \sqrt{7}$.
	
	Diện tích tam giác $A B C$ là: $S_{\Delta A B C}=\dfrac{1}{2} C A \cdot C B \cdot \sin \mathbb{C}=\dfrac{1}{2} A B \cdot C H$ $$\Rightarrow C H=\dfrac{C A \cdot C B \cdot \sin 120^{\circ}}{A B}=\dfrac{a \cdot 2 a \cdot \dfrac{\sqrt{3}}{2}}{a \sqrt{7}}=\dfrac{a \sqrt{21}}{7}.$$
	Vậy $\mathrm{d}(C C’,(A B B’ A’))=C H=\dfrac{a \sqrt{21}}{7}$.
	\itemch Ta có $A M$ và $C C’$ là hai đường thẳng chéo nhau mà $\heva{&C C’ / /(A B B’ A’) \\ &A M \subset(A B B’ A’)}$
	nên $\mathrm{d}(C C’, A M)=\mathrm{d}(C C’,(A B B’ A’))=\dfrac{a \sqrt{21}}{7}$.
	\itemch Vì $A B C.A’ B’ C’$ là hình lăng trụ đứng nên $A A’ \perp(A B C), A A’ \perp(A’ B’ C’)$. 
	\itemch Do vậy $\mathrm{d}((A B C),(A’ B’ C’))=A A’=2 a$.\\
	Khối lăng trụ $A B C \cdot A’ B’ C’$ có chiều cao $h=A A’=2 a$.\\
	 Diện tích đáy là:
	$S=S_{A B C}=\dfrac{1}{2} C A \cdot C B \cdot \sin 120^{\circ}=\dfrac{1}{2} a \cdot 2 a \cdot \dfrac{\sqrt{3}}{2}=\dfrac{a^2 \sqrt{3}}{2}.$\\
	Thể tích khối lăng trụ là: $V=S h=\dfrac{a^2 \sqrt{3}}{2} \cdot 2 a=a^3 \sqrt{3}$.
\end{itemchoice}
}
\end{ex}

\Closesolutionfile{ans}
\indapan{3}{ans/ans-11-C8-B4-DS}

\caukq
\Opensolutionfile{ans}[ans/ans-11-C8-B4-KQ]
\begin{ex}%[1H8V5-5]
	Cho tứ diện đều ${A B C D}$ có cạnh ${4\sqrt{2}}$. Đoạn vuông góc chung của hai đường thẳng chéo nhau ${A B, C D}$ bằng
	\shortans{$4$}
\loigiai{
	\begin{center}
		\begin{tikzpicture}[scale=0.8, font=\footnotesize,line join=round, line cap=round, >=stealth]
			\path 
			(0,0) coordinate (B)
			++(-60:2) coordinate (C)
			(6,0) coordinate (D)
			($(D)!1/2!(C)$) coordinate (J)
			($(B)!2/3!(J)$) coordinate (H)
			($(H)+(0,4)$) coordinate (A)
			($(A)!1/2!(B)$) coordinate (I)
			;
			\foreach \i in{B,C,D}{\draw (A)--(\i);};
			\draw (B)--(C)--(D) (A)--(J) (I)--(C);
			\draw[dashed] (B)--(D) (B)--(J) (D)--(I)--(J);
			\pic[draw,angle eccentricity=1.8,angle radius=2mm]{right angle=B--J--C};
			\pic[draw,angle eccentricity=1.8,angle radius=2mm]{right angle=B--I--C};
			\pic[draw,angle eccentricity=1.8,angle radius=2mm]{right angle=D--I--A};
			\pic[draw,angle eccentricity=1.8,angle radius=2mm]{right angle=A--J--D};
			\foreach \i/\g in {A/90,B/180,C/-90,D/0,J/-90}
			\fill[black] (\i) circle(1pt)+(\g:4mm)node[scale=1]{$\i$};
		\end{tikzpicture}
	\end{center}
	Gọi ${I, J}$ theo thứ tự là trung điểm của ${A B, C D}$.\\
	Các tam giác ${A B C, A B D}$ đều có ${I}$ là trung điểm ${A B}$ nên
	$\heva{&A B \perp C I \\
			&A B \perp D I
	} \Rightarrow A B \perp(I C D)$.
	Mà ${I J \subset(I C D) \Rightarrow A B \perp I J}$. \tagEX{1}
	\noindent Tương tự, các tam giác ${A C D, B C D}$ đều có ${J}$ là trung điểm ${C D}$ nên\\
	$\heva{&
			CD\perp AJ \\
			&CD\perp BJ
		}\Rightarrow CD\perp (ABJ)$.
		Mà $J\subset (JAB)\Rightarrow CD\perp IJ$\tagEX{2}
	\noindent Từ (1) và (2) suy ra ${I J}$ là đoạn vuông góc chung của hai đường thẳng ${A B, C D}$.\\
	Ta có ${C I=\dfrac{2 a \cdot \sqrt{3}}{2}=a \sqrt{3};  I J=\sqrt{C I^2-C J^2}=\sqrt{3 a^2-a^2}=a \sqrt{2}}$.\\
	Thay $a=\sqrt{8}$ thì  d (${A B, C D}$)= $4$.
}
\end{ex}
\begin{ex}%[1H8V7-5]
\immini
{
	Một cái hộp hình lập phương, bên trong nó đựng một mô hình đồ chơi có dạng hình chóp tứ giác đều mà đỉnh của hình chóp đó trùng với tâm của một mặt chiếc hộp, giả sử hình vuông đáy của hình chóp trùng với một mặt của chiếc hộp (mặt này cùng với mặt chứa đỉnh hình chóp là hai mặt đối nhau). Biết cạnh của chiếc hộp bằng ${30 {~cm}}$, hãy tính thể tích phần không gian bên trong chiếc hộp không bị chiếm bởi mô hình đồ chơi dạng hình chóp (mô hình đồ chơi được làm bởi chất liệu nhựa đặc bên trong, kết quả có đơn vị là dm$^3$).
}{
	\begin{tikzpicture}[scale=1, font=\footnotesize,line join=round, line cap=round, >=stealth]
		\path 
		(0,0) coordinate (A)
		++(-130:2) coordinate (B)
		++(0:4) coordinate (C)
		($(A)+(C)-(B)$) coordinate (D)
		
		;
		\foreach \i in {A,B,C,D}{
			\coordinate (\i') at ($(\i)+(0,4)$);
		}
		\coordinate (O) at ($(A')!1/2!(C')$);
		\draw (A')--(B')--(C')--(D')--cycle;
		\draw (B)--(B') (C)--(C') (D)--(D')  (B)--(C)--(D);
		\draw[dashed,thin](B)--(A)--(A') (A)--(D) (O)--(A) (O)--(B) (O)--(C) (O)--(D);
		\foreach \i/\g in {A'/90,B'/90,C'/90,D'/90,A/-90,B/-90,C/-90,D/-90}
		\fill[black] (\i) circle(1pt)+(\g:5mm);
	\end{tikzpicture}
}
\shortans{$18$}
\loigiai{
Thể tích cái hộp (khối lập phương) là: $V_1=30^3=27000$ (cm$^3$).\\
Xét đồ chơi có dạng hình chóp tứ giác đều, chiều cao của hình chóp bằng với một cạnh của hình lập phương, hay $h=30$ cm, đáy của hình chóp có diện tích $S=30^2=900$ (cm$^2$).\\
Thể tích khối đồ chơi (khối chóp tứ giác đều) là:
$V_2=\dfrac{1}{3} S h=\dfrac{1}{3} \cdot 900 \cdot 30=9000$ (cm$^3$).\\
Thể tích phần không gian bên trong chiếc hộp không bị chiếm bởi mô hình đồ chơi dạng hình chóp: $V=V_1-V_2=27000-9000=18000$ (cm$^3$)=18 (dm$^3$).
}
\end{ex}
\begin{ex}%[1H8V7-6]
Một hình chóp cụt đều ${A B C \cdot A’ B’ C’}$ có cạnh đáy lớn bằng ${4 a}$, cạnh đáy nhỏ bằng ${2 a}$ và chiều cao của nó bằng ${\dfrac{3 a}{2}}$. Tìm thể tích của khối chóp cụt đều đó với $a=2\sqrt{3}$.
\shortans{$252$}
\loigiai{
	\begin{center}
		\begin{tikzpicture}[scale=1, font=\footnotesize,line join=round, line cap=round, >=stealth]
			\path 
			(0,0) coordinate (A)
			++(-60:3) coordinate (B)
			(6,0) coordinate (C)
			($(B)!1/2!(C)$) coordinate (K)
			($(A)!2/3!(K)$) coordinate (O)
			($(O)+(0,5)$) coordinate (S)
			($(S)!.5!(A)$) coordinate (A')
			($(S)!.5!(B)$) coordinate (B')
			($(S)!.5!(C)$) coordinate (C')
			($(B')!1/2!(C')$) coordinate (J)
			($(A')!2/3!(J)$) coordinate (I)
			($(A)!.3!(K)$) coordinate (H)
			;
			\foreach \i in{A,B,C}{\draw (\i')--(\i);};
			\draw (A)--(B)--(C) (A')--(B')--(C')--cycle (A')--(J);
			\draw[dashed] (I)--(O) (A)--(C) (A)--(K) (A')--(H);
			\pic[draw,angle eccentricity=1.8,angle radius=2mm]{right angle=S--O--K};
			\pic[draw,angle eccentricity=1.8,angle radius=2mm]{right angle=A'--H--O};
			\pic[draw,angle eccentricity=1.8,angle radius=2mm]{right angle=O--I--J};
			\foreach \i/\g in {A/180,B/-120,C/0,K/-50,O/-90,A'/90,B'/180,C'/90,J/-50,I/30,H/-90}
			\fill[black] (\i) circle(1pt)+(\g:4mm)node[scale=1]{$\i$};
		\end{tikzpicture}
	\end{center}
Gọi ${O, I}$ theo thứ tự là tâm của đáy lớn ${A B C}$ và đáy bé ${A’ B’ C’;  K, J}$ theo thứ tự là trung điểm của ${B C}$ và ${B’ C’}$.\\
Ta có ${h=I O=\dfrac{3 a}{2}}$ là chiều cao của hình chóp cụt đều ${A B C \cdot A’ B’ C’}$.\\
Diện tích hai đáy hình chóp cụt đều là\\
${S_1=S_{\Delta A B C}=\dfrac{(4 a)^2 \sqrt{3}}{4}=4 a^2 \sqrt{3};  S_2=S_{\Delta A’ B’ C’}=\dfrac{(2 a)^2 \sqrt{3}}{4}=a^2 \sqrt{3}
}$.\\
Thể tích khối chóp cụt đều là\\
$V=\dfrac{1}{3}h\left(S_1+\sqrt{S_1S_2}+S_2\right)$
$=\dfrac{1}{3}\cdot \dfrac{3a}{2}\left(4a^2\sqrt{3}+\sqrt{4a^2\sqrt{3}\cdot a^2\sqrt{3}}+a^2\sqrt{3}\right)=\dfrac{7a^3\sqrt{3}}{2}$ (đơn vị thể tích).\\
Với $a=2\sqrt{3}$ thì $V=252$.
}
\end{ex}
\begin{ex}%[1H8V5-2]
Cho hình chóp ${S.ABCD}$ có ${SA \perp(ABCD), SA=10, ABCD}$ là hình vuông cạnh bằng $5$. Gọi ${O}$ là tâm của ${A B C D}$. Khoảng cách từ $S$ đến $DM$ với $M$ là trung điểm ${O C}$ và làm tròn kết quả đến một chữ số thập phân.
\shortans{$13{,}8$}
\loigiai{
\begin{center}
	\begin{tikzpicture}[scale=1, font=\footnotesize,line join=round, line cap=round, >=stealth]
		\path 
		(0,0) coordinate (A)
		++(-140:3) coordinate (D)
		++(0:5) coordinate (C)
		($(A)+(C)-(D)$) coordinate (B)
		($(A)+(0,4)$) coordinate (S)
		($(A)!.5!(C)$) coordinate (O)
		($(O)!.5!(C)$) coordinate (M)
		($(D)!(S)!(M)$) coordinate (K)
		($(S)!(A)!(C)$) coordinate (H)
		;
		\foreach \i in{B,C,D}{\draw (S)--(\i);};
		\draw (B)--(C)--(D);
		\draw[dashed] (S)--(A)--(B) (A)--(D) (A)--(C) (B)--(D)--(M) (S)--(K)--(A)--(H);
		\pic[draw,angle eccentricity=1.8,angle radius=2mm]{right angle=A--H--C};
		\pic[draw,angle eccentricity=1.8,angle radius=2mm]{right angle=S--K--M};
		\path (O)--(M) node[midway,sloped]{$/$};
		\path (C)--(M) node[midway,sloped]{$/$};
		\path (S)--(A) node[left,midway]{$2a$};
		\path (D)--(C) node[below,midway]{$a$};
		\foreach \i/\g in {A/180,B/-90,C/-90,D/-90,S/90,O/90,K/-50,H/0,M/50}
		\fill[black] (\i) circle(1pt)+(\g:3mm)node[scale=1]{$\i$};
	\end{tikzpicture}
\end{center}
Kẻ ${S K \perp D M}$ tại ${K \Rightarrow \mathrm{d}(S, D M)=S K}$.\\
Ta có ${\heva{&D M \perp S A \\ &D M \perp S K} \Rightarrow DM \perp(SAK) \Rightarrow DM \perp A K}$\\
Ta có ${\Delta K M A \backsim \Delta O M D}$.\\
${\Rightarrow \dfrac{K A}{O D}=\dfrac{A M}{D M} \Rightarrow K A=\dfrac{A M \cdot O D}{D M}=\dfrac{\dfrac{3}{4} a \sqrt{2} \cdot a \sqrt{2}}{\sqrt{\left(\dfrac{a \sqrt{2}}{2}\right)^2+\left(\dfrac{a \sqrt{2}}{4}\right)^2}}=\dfrac{3 \sqrt{10}}{5} a}$.\\
Ta có ${S K=\sqrt{S A^2+A K^2}=\sqrt{(2 a)^2+\left(\dfrac{3 \sqrt{10}}{5} a\right)^2}=\dfrac{\sqrt{190}}{5} a}$\\
Vậy ${\mathrm{d}(S, D M)=\dfrac{\sqrt{190}}{5} a}\approx 13{,}8$ với $a=5$.
}
\end{ex}
\begin{ex}%[1H8V5-3]
Cho hình chóp ${S.  A B C D}$ có đáy là hình vuông cạnh $6$, hai mặt phẳng ${(S A B)}$ và ${(S B C)}$ cùng vuông góc với mặt phẳng ${(A B C D)}$ và ${S C=6\sqrt{5}}$. Tính khoảng cách từ ${D}$ đến mặt phẳng ${(S A C)}$.
\shortans{$4$}
\loigiai{
		\begin{center}
			\begin{tikzpicture}[scale=1, font=\footnotesize,line join=round, line cap=round, >=stealth]
				\path 
				(0,0) coordinate (B)
				++(-140:3) coordinate (A)
				++(0:5) coordinate (D)
				($(B)+(D)-(A)$) coordinate (C)
				($(B)+(0,4)$) coordinate (S)
				($(A)!.5!(C)$) coordinate (O)
				($(S)!(B)!(O)$) coordinate (H)
				;
				\foreach \i in{A,C,D}{\draw (S)--(\i);};
				\draw (A)--(D)--(C);
				\draw[dashed] (S)--(B)--(A) (A)--(C)--(B)--(D) (S)--(O) (B)--(H);
				\pic[draw,angle eccentricity=1.8,angle radius=2mm]{right angle=B--H--O};
				\path (S)--(C) node[right,midway]{$a\sqrt{5}$};
				\path (D)--(A) node[below,midway]{$a$};
				\foreach \i/\g in {A/180,B/180,C/-90,D/-90,S/90,O/-90,H/0}
				\fill[black] (\i) circle(1pt)+(\g:3mm)node[scale=1]{$\i$};
			\end{tikzpicture}
		\end{center}
	Đặt $AB=a, SC=a\sqrt{5}$.\\
	Kẻ $BH\perp SO$ tại ${H}$.\\
	Ta có  $\heva{&A C \perp S B \\ &A C \perp B D} \Rightarrow AC \perp(SBD) \Rightarrow AC \perp BH$.\\
	Ta lại có  ${B H \perp S O \Rightarrow B H \perp(S A C) \Rightarrow \mathrm{d}(B,(S A C))=B H}$.\\
	Ta có  ${S B=\sqrt{S C^2-B C^2}=\sqrt{(a \sqrt{5})^2-a^2}=2 a}$.\\
	Ta có  ${B H=\dfrac{1}{\sqrt{\dfrac{1}{S B^2}+\dfrac{1}{O B^2}}}=\dfrac{1}{\sqrt{\dfrac{1}{4a^2}+\dfrac{1}{\frac{a^2}{2}}}}}=\dfrac{2a}{3}$. \\
	Vậy ${\mathrm{d}(B,(S A C))=\dfrac{2}{3} a}$.\\
	Ta có  ${D B}$ cắt ${(S A C)}$ tại ${O}$
	$\Rightarrow \dfrac{\mathrm{d}(D,(S A C))}{\mathrm{d}(B,(S A C))}=\dfrac{D O}{B O}=1$\\
	$ \Rightarrow \mathrm{d}(D,(S A C))=\mathrm{d}(B,(S A C))=\dfrac{2}{3} a=4$ với $a=6$.
}
\end{ex}
\begin{ex}%[1H8V5-3]
	Cho hình chóp ${S.  A B C}$ có đáy là tam giác vuông cân tại ${A, S H \perp(A B C)}$ với ${H}$ là trung điểm ${B C}$. Biết ${A B=S C=\sqrt{6}}$. Tính khoảng cách từ ${B}$ đến mặt phẳng ${(S A C)}$.
	\shortans{$2$}
	\loigiai{
		\begin{center}
			\begin{tikzpicture}[scale=1, font=\footnotesize,line join=round, line cap=round, >=stealth]
				\path
				(0,0) coordinate (B)
				++(-45:2) coordinate (A)
				(5,0) coordinate (C)
				($(B)!1/2!(C)$) coordinate (H)
				($(A)!1/2!(C)$) coordinate (I)
				($(H)+(0,4)$) coordinate (S)
				($(S)!1/2!(I)$) coordinate (K)
				;
				\foreach \i in {A,B,C}{\draw (S)--(\i);}
				\draw (B)--(A)--(C) (S)--(I);
				\draw[dashed,thin](B)--(C) (S)--(H)--(I) (H)--(K);
				\pic[draw,thin,angle radius=2mm] {right angle = S--H--C};
				\pic[draw,thin,angle radius=2mm] {right angle = B--A--C};
				\path (A)--(B) node[midway,left]{$a$};
				\path (S)--(C) node[midway,right]{$a$};
				\path (B)--(A) node[midway,sloped,rotate=90]{$=$};
				\path (I)--(A) node[midway,sloped,rotate=90]{$=$};
				\path (B)--(H) node[midway,sloped]{$|$};				
				\path (C)--(H) node[midway,sloped]{$|$};
				\foreach \i/\g in {S/90,A/-90,B/180,C/0,H/-90,I/-50,K/0}
				\fill[black] (\i) circle(1pt)+(\g:4mm)node[scale=1]{$\i$};
			\end{tikzpicture}
		\end{center}
	Kẻ ${H I \perp A C}$, kẻ ${H K \perp S I}$ tại ${K}$.\\
	Ta có  ${\left\{\begin{array}{l}A C \perp S H \\ A C \perp H I\end{array} \Rightarrow A C \perp(S H I) \Rightarrow A C \perp H K\right.}$\\
	Ta lại có  ${H K \perp S I \Rightarrow H K \perp(S A C) \Rightarrow \mathrm{d}(H,(S A C))=H K}.$\\
	Ta có  ${{HI}=\dfrac{1}{2} a}$;
	${S H=\sqrt{S C^2-H C^2}=\sqrt{a^2-\left(\dfrac{a \sqrt{2}}{2}\right)^2}=\dfrac{\sqrt{2}}{2} a}$.\\
	Ta có  ${H K=\dfrac{1}{\sqrt{\dfrac{1}{S H^2}+\dfrac{1}{H I^2}}}=\dfrac{1}{\sqrt{\dfrac{1}{\sqrt{2})^2}+\dfrac{1}{(1)^2}}}=\dfrac{\sqrt{6}}{6} a}$.\\
	Vậy ${\mathrm{d}(H,(S A C))=\dfrac{\sqrt{6}}{6} a}$.\\
	Ta có  ${B H}$ cắt ${(S A C)}$ tại ${C}$
	$\Rightarrow \dfrac{\mathrm{d}(B,(S A C))}{\mathrm{d}(H,(S A C))}=\dfrac{B C}{H C}=2$.\\
$ \Rightarrow \mathrm{d}(B,(S A C))=2 \mathrm{d}(H,(S A C))=\dfrac{\sqrt{6}}{3} a=2$ với $a=\sqrt{6}$.
	}
\end{ex}

\Closesolutionfile{ans}
\indapan{6}{ans/ans-11-C8-B4-KQ}


\subsection{Đề 2}
 \Opensolutionfile{ans}[ans/ans-De1]
 \caulc
 \begin{ex}%[1H8H5-3] 
 	Cho hình chóp $S.ABC$ có đáy $ABC$ là tam giác cân tại $B$, cạnh bên $SA$ vuông góc với đáy, $J$ là hình chiếu của $A$ lên $BC$. Kí hiệu $\mathrm{d}[A,(SBC)]$ là khoảng cách giữa điểm $A$ đến mặt phẳng $(SBC)$. Khẳng định nào sau đây đúng?
 	\choice
 	{$\mathrm{d}[A,(SBC)]=AK$ với $K$ là hình chiếu của $A$ lên $SB$}
 	{\True $\mathrm{d}[A,(SBC)]=AK$ với $K$ là hình chiếu của $A$ lên $SJ$}
 	{$\mathrm{d}[A,(SBC)]=AK$ với $K$ là hình chiếu của $A$ lên $SC$}
 	{$\mathrm{d}[A,(SBC)]=AK$ với $K$là hình chiếu của $A$ lên $SM$}
 	\loigiai{
 		\immini{
 			Ta có $\heva{&SA\perp BC\\&AJ\perp BC}\Rightarrow BC\perp (SAJ)$.\\
 			Gọi $K$ là hình chiếu của $A$ lên $SJ$ ta có $$\heva{&AK\perp BC\,(BC\perp (SAJ))\\&AJ\perp BC.}$$
 			Suy ra $AK\perp (SBC)$.\\
 			Do đó $\mathrm{d}[A,(SBC)]=AK$.
 		}
 		{
 			\begin{tikzpicture}[line join=round, line cap=round,thick]
 				\coordinate (A) at (-2,0);
 				\coordinate (B) at (0,-3);
 				\coordinate (C) at (3,0);
 				\coordinate (S) at ($(A)+(0,4)$);
 				\coordinate (J) at ($(B)!0.4!(C)$);
 				\coordinate (K) at ($(S)!0.5!(J)$);
 				\draw(S)--(A) (S)--(B) (S)--(C) (B)--(C) (A)--(B)(S)--(J);
 				\draw[dashed,thin](A)--(C)(K)--(A)--(J);
 				\pic[draw,thin,angle radius=2mm] {right angle = S--A--B} 
 				pic[draw,thin,angle radius=2mm] {right angle = S--A--C}
 				pic[draw,thin,angle radius=2mm] {right angle = A--J--B}
 				pic[draw,thin,angle radius=2mm] {right angle = A--K--J};
 				\foreach \i/\g in {S/90,A/180,B/-90,C/0,J/0,K/0}{\draw[fill=white](\i) circle (1.5pt) ($(\i)+(\g:3mm)$) node[scale=1]{$\i$};}
 			\end{tikzpicture}
 		}
 	}
 \end{ex}
 
 \begin{ex}%[1H8N7-6]
 	Cho khối hộp chữ nhật có 3 kích thước $3;\,4;\,5$. Thể tích của khối hộp đã cho bằng?
 	\choice
 	{$10$}
 	{$20$}
 	{$12$}
 	{\True $60$}
 	\loigiai{
	Thể tích khối hộp chữ nhật $V=a\cdot b \cdot c=3\cdot 4\cdot 5=60$ (đvtt).
 	}
 \end{ex}
 
 \begin{ex}%[1H8H7-3]
 	Cho khối chóp $S.ABCD$ có đáy $ABCD$ là hình vuông cạnh $a$, tam giác $SAB$ cân tại $S$ và nằm trong mặt phẳng vuông góc với đáy, $SA=2a$. Tính theo $a$ thể tích khối chóp $S.ABCD$
 	\choice
 	{$V=2a^3$}
 	{$V=\dfrac{a^3\sqrt{15}}{12}$}
 	{\True $V=\dfrac{a^3\sqrt{15}}{6}$}
 	{$V=\dfrac{2a^3}{3}$}
 	\loigiai{
 		\begin{center}
 			\begin{tikzpicture}[scale=0.7, font=\footnotesize, line join=round, line cap=round, >=stealth]
 			
 			\path
 			(0,0) coordinate (A)
 			($(A)+(-150:3)$) coordinate (B)
 			($(A)+(0:4)$) coordinate (D)
 			($(B)+(D)-(A)$) coordinate (C)
 			($(B)!0.5!(A)$) coordinate (H)
 			($(H)+(90:5)$) coordinate (S)
 			;
 			
 			\draw
 			(S)--(B)--(C)--(D)--(S)--(C)
 			;
 			
 			\draw[dashed]
 			(B)--(A)--(D)
 			(H)--(S)--(A)
 			;
 			
 			\foreach \x/\g in {A/160,B/-135,C/-90,D/0,H/-10,S/90}
 			\fill[black] (\x) circle (0.6pt)+(\g:.4)node{$\x$};
 			\pic[draw,thin,angle radius=2mm] {right angle = S--H--A};
 		\end{tikzpicture}	
 		\end{center}
 		Kẻ $SH \perp AB$ tại $H$, suy ra $SH=\sqrt{SA^2-AH^2}=\sqrt{4a^2-\dfrac{1}{4}a^2}=\dfrac{\sqrt{15}}{2}a$.\\
 		Ta có $ (SAB) $ vuông góc với mặt phẳng đáy $ (ABCD) $ do đó chiều cao của hình chóp là $SH$ cũng chính là chiều cao của $ \triangle SAB $.\\
 		Thể tích khối chóp $V=\dfrac{1}{3}\cdot SH\cdot S_{ABCD}=\dfrac{1}{3}\cdot\dfrac{\sqrt{15}}{2}a \cdot a^2=\dfrac{\sqrt{15}}{6}a^3$ (đvtt).
 	}
 \end{ex}
\begin{ex}%[1H8H5-4] 
	Cho hình chóp $S.ABCD$ có đáy là hình vuông cạnh $a$, $SA$ vuông góc với đáy, $SA=a$. Khoảng cách giữa hai đường thẳng $SB$ và $CD$ là 
	\choice
	{\True $a$}
	{$2a$}
	{$a\sqrt{2}$}
	{$a\sqrt{3}$}
	\loigiai{
	\begin{center}
		\begin{tikzpicture}[scale=0.7, line join=round, line cap=round,thick]
			\coordinate (A) at (0,0);
			\coordinate (B) at (-2,-3);
			\coordinate (D) at (7,0);
			\coordinate (C) at ($(B)+(D)-(A)$);
			\coordinate (S) at ($(A)+(0,5)$);
			\draw(S)--(B) (S)--(C) (S)--(D) (B)--(C)--(D);
			\draw[dashed,thin](S)--(A) (A)--(B) (A)--(D);
			\pic[draw,thin,angle radius=3mm] {right angle = S--A--D} pic[draw,thin,angle radius=3mm] {right angle = S--A--B};
			\foreach \i/\g in {S/90,A/-90,B/-90,C/-90,D/0}{\draw[fill=white](\i) circle (1.5pt) ($(\i)+(\g:3mm)$) node[scale=1]{$\i$};}
		\end{tikzpicture}
	\end{center}
	Ta có $\heva{&SA\perp BC\\&AB\perp BC}\Rightarrow BC\perp (SAB)\Rightarrow BC\perp SB \subset (SAB)$.\\
	Mà $BC\perp CD$ ($ABCD$ là hình vuông), suy ra $\mathrm{d}[SB,CD]=BC=a$.
	}
\end{ex}
\begin{ex}%[1H8H5-3] 
	Cho hình chóp $S.ABC$ có $ABC$ tam giác đều cạnh $2a$, $SA\perp \left( ABC \right)$ và $SA=a\sqrt{6}$. Gọi $M$ là trung điểm của $BC$, khi đó khoảng cách từ $A$ đến đường thẳng $SM$ bằng
	\choice
	{\True $a\sqrt{2}$}
	{$a\sqrt{3}$}
	{$a\sqrt{6}$}
	{$a\sqrt{11}$}
	\loigiai{
	\begin{center}
		\begin{tikzpicture}[line join=round, line cap=round,thick]
			\coordinate (A) at (-2,0);
			\coordinate (B) at (0,-3);
			\coordinate (C) at (3,0);
			\coordinate (S) at ($(A)+(0,4)$);
			\coordinate (M) at ($(B)!0.5!(C)$);
			\coordinate (K) at ($(S)!0.5!(M)$);
			\draw (M)--(S)--(A) (S)--(B) (S)--(C) (B)--(C) (A)--(B);
			\draw[dashed,thin](M)--(A)--(C)(A)--(K);
			\pic[draw,thin,angle radius=2mm] {right angle = S--A--B} 
			pic[draw,thin,angle radius=2mm] {right angle = A--K--M}
			pic[draw,thin,angle radius=2mm] {right angle = S--A--C};
			\foreach \i/\g in {S/90,A/180,B/-90,C/0,M/0,K/0}{\draw[fill=white](\i) circle (1.5pt) ($(\i)+(\g:3mm)$) node[scale=1]{$\i$};}
		\end{tikzpicture}
	\end{center}
	Ta có tam giác $ABC$ đều và $M$ là trung điểm $BC$, suy ra $AM= a\sqrt{3}$.\\
	Kẻ $AK \perp SM$, suy ra $\mathrm{d}[A,SM]=AK=\sqrt{\dfrac{SA^2\cdot AM^2}{SA^2+AM^2}}=\sqrt{\dfrac{6a^2\cdot 3a^2}{6a^2 + 3a^2}}=a\sqrt{2}$.
	}
\end{ex}
\begin{ex}%[1H8H5-3]
	Cho hình hộp chữ nhật $ABCD.A' B' C' D' $ có $AB=a$, $AD=\sqrt{3}a$, $AA' =2\sqrt{3}a$ (tham khảo hình vẽ). Khoảng cách giữa hai mặt phẳng $(ABCD)$ và $(A'B'C'D')$ bằng
	\choice
	{$a$}
	{\True $2a\sqrt{3}$}
	{$3a$}
	{ $a\sqrt{3}$}
	\loigiai{
		\immini{
		Vì $(ABCD)\parallel (A'B'C'D')$ nên $$\mathrm{d}[(ABCD),(A'B'C'D')]=\mathrm{d}[A,(A'B'C'D')]=AA'=2a\sqrt{3}.$$
		}
		{
			\begin{tikzpicture}[font=\footnotesize, line join=round, line cap=round, >=stealth]
				\coordinate  (A) at (0,0); 
				\coordinate  (A') at (0,2);
				\coordinate (B) at (2,0);
				\coordinate  (B') at (2,2);
				\coordinate (C) at (3,0.5);
				\coordinate  (C') at (3,2.5);
				\coordinate (D) at (1,0.5);
				\coordinate  (D') at (1,2.5);
				\draw (A) -- (B) -- (C) -- (C') -- (D') -- (A') -- (B') -- (C') (A) -- (A') (B) -- (B');
				\draw[dashed] (A)--(D) -- (C) (D) --(D'); 
				\foreach \i/\g in {A/180,B/-90,C/0,D/160,A'/180,B'/0,C'/90,D'/90}{\draw[fill=white](\i) circle (1.5pt) ($(\i)+(\g:3mm)$) node[scale=1]{$\i$};}
			\end{tikzpicture}
		}
		
		
	}
\end{ex}
\begin{ex}%[1H8H5-4] 
	\immini{
	Cho hình hộp chữ nhật $ABCD.A'B'C'D'$ có $AB=1,BC=2;AA'=3$ (tham khảo hình vẽ).
	Khoảng cách giữa hai đường $AB'$ và $CC'$ bằng?
	\choice
	{$1$}
	{\True $2$}
	{$3$}
	{$\sqrt{2}$}
	}
	{
	\begin{tikzpicture}[font=\footnotesize, line join=round, line cap=round, >=stealth]
		\coordinate  (A) at (0,0); 
		\coordinate  (A') at (0,2);
		\coordinate (B) at (2,0);
		\coordinate  (B') at (2,2);
		\coordinate (C) at (3,0.5);
		\coordinate  (C') at (3,2.5);
		\coordinate (D) at (1,0.5);
		\coordinate  (D') at (1,2.5);
		\draw (A) -- (B) -- (C) -- (C') -- (D') -- (A') -- (B') -- (C') (A) -- (A') (B) -- (B');
		\draw[dashed] (A)--(D) -- (C) (D) --(D'); 
		\foreach \i/\g in {A/180,B/-90,C/0,D/160,A'/180,B'/0,C'/90,D'/90}{\draw[fill=white](\i) circle (1.5pt) ($(\i)+(\g:3mm)$) node[scale=1]{$\i$};}
	\end{tikzpicture}
	}

	\loigiai{
		Vì $\heva{& AB'\subset (ABB'A') \\
				& CC'\subset (CDD'C') \\ 
			 	& (ABB'A')\parallel(CDD'C')}.$\\
		nên $$\mathrm{d}\left[AB',CC'\right]=\mathrm{d}\left[\left(ABB'A'\right),\left(CDD'C'\right)\right]=\mathrm{d}\left[(A,\left(CDD'C'\right))\right]=AD=2.$$
	}
\end{ex}
\begin{ex}%[1H8N5-3]
Cho hình chóp $S.ABCD$ có đáy là hình chữ nhật với $AB=a$, $AD=3a$. Biết $SA=2a$ và $SA$ vuông góc với mặt phẳng $(ABCD)$. Khoảng cách từ điểm $C$ đến đường thẳng $AD$ bằng
\choice
{\True $a$}
{$2a$}
{$3a$}
{$a\sqrt{10}$}
\loigiai{
	\begin{center}
		\begin{tikzpicture}[line join=round, line cap=round,thick]
		\coordinate (A) at (0,0);
		\coordinate (B) at (-2,-3);
		\coordinate (D) at (7,0);
		\coordinate (C) at ($(B)+(D)-(A)$);
		\coordinate (S) at ($(A)+(0,5)$);
		\draw(S)--(B) (S)--(C) (S)--(D) (B)--(C)--(D);
		\draw[dashed,thin](S)--(A) (A)--(B) (A)--(D);
		\pic[draw,thin,angle radius=3mm] {right angle = S--A--D} pic[draw,thin,angle radius=3mm] {right angle = S--A--B};
		\foreach \i/\g in {S/90,A/-90,B/-90,C/-90,D/0}{\draw[fill=white](\i) circle (1.5pt) ($(\i)+(\g:3mm)$) node[scale=1]{$\i$};}
	\end{tikzpicture}
	\end{center}
	Ta có $CD\perp AD$ nên $\mathrm{d}(C,AD)=CD=a$.
}
\end{ex}

\begin{ex}%[1H4K7-4]
	Cho khối lăng trụ đều $ABC.A'B'C'$ có cạnh đáy bằng $a$. Khoảng cách từ điểm $A'$ đến mặt phẳng $(AB'C')$ bằng $\dfrac{2a \sqrt{3}}{\sqrt{19}}$. Thể tích của khối lăng trụ đã cho là
	\choice
	{$\dfrac{a^3\sqrt{3}}{4}$}
	{$\dfrac{a^3\sqrt{3}}{6}$}
	{\True $\dfrac{a^3\sqrt{3}}{2}$}
	{$\dfrac{3a^3}{2}$}
	\loigiai{
		\begin{center}
			\begin{tikzpicture}[scale=1, font=\footnotesize,line join=round, line cap=round, >=stealth]
				\def\b{5.5}  % Chiều dài BC
				\def\c{1.8}  % Chiều dài AB
				\def\gA{-40} % Góc A
				\def\h{3.5}  % Chiều cao hình chóp
				\path  (0,0) coordinate (A)
				(0:\b) coordinate (C)
				(\gA:\c) coordinate (B)
				(A)+(-90:\h) coordinate (A')
				(A')+(0:\b) coordinate (C')
				(A')+(\gA:\c) coordinate (B')
				($(B')!0.5!(C')$) coordinate (M)
				($(A)!(A')!(M)$) coordinate (H)	% ($(B)!(A)!(C)$) Hình chiếu của A lên BC
				
				;
				\draw (A')--(B')--(C') (A)--(B)--(C)--cycle (A)--(A') (B)--(B') (C)--(C');
				\draw[dashed] (C')--(A)--(M)--(A')--(H) (A')--(C'); 
				\pic[draw,thin,angle radius=2mm] {right angle = C--A--A'};
				\pic[draw,thin,angle radius=2mm] {right angle = B--A--A'};
				\pic[draw,thin,angle radius=2mm] {right angle = A'--H--M};
				\foreach \i/\g in {A/135,B/90,C/45,A'/-135,B'/-45,C'/-45,M/-45,H/45}{\draw[fill=black](\i) circle (1pt) ($(\i)+(\g:3mm)$) node[scale=1]{$\i$};}
			\end{tikzpicture}
		\end{center}
		Gọi $M$ là trung điểm của $B'C'$.\\
		Ta có $\heva{&AA'\perp B'C' \\&A'M\perp B'C'} \Rightarrow B'C' \perp(AA'M) \Rightarrow(AB'C') \perp (AA'M)$ theo giao tuyến $AM$.\\
		Kẻ $A'H\perp AM$ trong mặt phẳng $(AA'M)$, suy ra $A'H\perp(AB'C')$.\\
		Vậy khoảng cách từ $A'$ đến mặt phẳng $(AB'C')$ là $A'H=\dfrac{2a \sqrt{3}}{\sqrt{19}}$.\\
		Ta có $\dfrac{1}{A'H^2}=\dfrac{1}{A'A^2}+\dfrac{1}{A'M^2} \Rightarrow \dfrac{1}{A'A^2}=\dfrac{1}{A'H^2}-\dfrac{1}{A'M^2}=\dfrac{1}{4a^2} \Rightarrow A'A=2a$.\\
		Vậy thể tích khối lăng trụ là $V=AA'\cdot S_{A'B'C'}=2a \cdot \dfrac{a^2\sqrt{3}}{4}=\dfrac{a^3\cdot \sqrt{3}}{2}$. 
	}
\end{ex}
\begin{ex}%[1H4K5-4]
	Cho tứ diện $ABCD$ đều có cạnh bằng $2\sqrt{2}$. Gọi $G$ là trọng tâm tứ diện $ABCD$ và $M$ là trung điểm $AB$. Khoảng cách giữa hai đường thẳng $BG$ và $CM$ bằng
	\choice
	{\True $\dfrac{2}{\sqrt{14}}$}
	{$\dfrac{2}{\sqrt{5}}$}
	{$\dfrac{3}{2\sqrt{5}}$}
	{$\dfrac{2}{\sqrt{10}}$}
	\loigiai{
		\begin{center}
			\begin{tikzpicture}[scale=1, font=\footnotesize,line join=round, line cap=round, >=stealth]
				\def\b{7}  % Chiều dài AC
				\def\c{5}  % Chiều dài AB
				\def\gA{-22} % Góc A
				\def\h{4}  % Chiều cao hình chóp
				\path  (0,0) coordinate (B)
				(0:\b) coordinate (D)
				(\gA:\c) coordinate (C)
				($(D)!0.5!(C)$) coordinate (N)
				($(C)!0.5!(N)$) coordinate (K)
				($(B)!2/3!(N)$)			coordinate(H)
				($(H)+(78:\h)$) coordinate (A)
				($(A)!0.5!(B)$) coordinate (M)	
				($(M)!0.5!(N)$) coordinate (G)
				($(B)!0.55!(K)$) coordinate (I)
				($(I)!(H)!(G)$) coordinate (J)
				;
				\foreach \i in {D,B,C}{\draw (A)--(\i);}
				\draw (B)--(C)--(D) (M)--(C);
				\draw[dashed] (D)--(B)--(N) (G)--(B)--(K)--cycle (A)--(H) (M)--(N) (J)--(H)--(I)--(G);
				\pic[draw,thin,angle radius=2mm] {right angle = B--N--C};
				\pic[draw,thin,angle radius=2mm] {right angle = I--J--H};
				\pic[draw,thin,angle radius=2mm] {right angle = K--I--H};
				\foreach \i/\g in {A/90,B/-135,C/-45,D/-45,H/45,K/-45,I/-135,N/-45,M/135,G/40,J/135}{\draw[fill=black](\i) circle (1pt) ($(\i)+(\g:3mm)$) node[scale=1]{$\i$};}
			\end{tikzpicture}
		\end{center}
		Gọi $N$ là trung điểm $CD$, khi đó $G$ là trung điểm $MN$ và $AG$ đi qua trọng tâm $H$ của tam giác $BCD$.\\
		Ta có $AH\perp(BCD$. và $AH=\sqrt{AB^2-BH^2}=\sqrt{(2\sqrt{2})^2-\left(\dfrac{2\sqrt{6}}{3}\right)^2}=\dfrac{4\sqrt{3}}{3}$.\\
		Ta có $GH=\dfrac{1}{4} AH=\dfrac{\sqrt{3}}{3}$.\\
		Gọi $K$ là trung điểm $CN$ thì $GK\parallel CM$ nên $CM\parallel(BGK)$.\\
		Do đó
		$\mathrm{d}(BG; CM)=\mathrm{d}(C;(BGK))=\mathrm{d}(N;(BGK))=\dfrac{3}{2} \mathrm{d}(H;(BGK))$.\\
		Kẻ $HI\perp BK$, $HJ\perp GI$ với $I\in BK$, $J\in GI$.\\
		Khi đó $HJ\perp(BGK)$ và $HJ=\mathrm{d}(H;(BGK))$.\\
		Ta có $BK=\sqrt{BN^2+NK^2}=\sqrt{(\sqrt{6})^2+\left(\dfrac{\sqrt{2}}{2}\right)^2}=\dfrac{\sqrt{26}}{2}$.\\
		Ta có $HI=BH\cdot \sin \widehat{KBN}=BH\cdot \dfrac{KN}{BK}=\dfrac{2\sqrt{6}}{3} \cdot \dfrac{\dfrac{\sqrt{2}}{2}}{\dfrac{\sqrt{26}}{2}}=\dfrac{2\sqrt{6}}{3\sqrt{13}}$.\\
		Do đó $HJ=\dfrac{HI\cdot HG}{\sqrt{HI^2+HG^2}}=\dfrac{\dfrac{2\sqrt{6}}{3\sqrt{13}} \cdot \dfrac{\sqrt{3}}{3}}{\sqrt{\left(\dfrac{2\sqrt{6}}{3\sqrt{13}}\right)^2+\left(\dfrac{\sqrt{3}}{3}\right)^2}}=\dfrac{2\sqrt{2}}{3\sqrt{7}}$.\\
		Vậy $\mathrm{d}(BG; CM)=\dfrac{3}{2} \mathrm{d}(H;(BGK))=\dfrac{3}{2} HJ=\dfrac{3}{2} \cdot \dfrac{2\sqrt{2}}{3\sqrt{7}}=\dfrac{2}{\sqrt{14}}$. 
	}
\end{ex}

\begin{ex}%[1H4B5-3]
	Cho hình chóp $S.ABC$ có đáy là tam giác vuông cân tại $B$, $AB=a$. Biết $SA=2a$ và $SA$ vuông góc với mặt phẳng $(ABC)$. Khoảng cách từ điểm $A$ đến mặt phẳng $(SBC)$ bằng
	\choice
	{$a$}
	{$2a$}
	{\True $\dfrac{2a \sqrt{5}}{5}$}
	{$\dfrac{a \sqrt{5}}{2}$}
	\loigiai{
		\begin{center}
			\begin{tikzpicture}[scale=1, font=\footnotesize,line join=round, line cap=round, >=stealth]
				\def\b{4.5}  % Chiều dài BC
				\def\c{2}  % Chiều dài AB
				\def\gA{-40} % Góc A
				\def\h{3.5}  % Chiều cao hình chóp
				\path  (0,0) coordinate (A)
				(0:\b) coordinate (C)
				(\gA:\c) coordinate (B)
				(A)+(90:\h) coordinate (S)
				($(S)!(A)!(B)$) coordinate (H)
				;
				\foreach \i in {A,B,C}{\draw (S)--(\i);}
				\draw (H)--(A)--(B)--(C); 
				\draw[dashed] (A)--(C);
				\pic[draw,thin,angle radius=2mm] {right angle = C--A--S};
				\pic[draw,thin,angle radius=2mm] {right angle = B--A--S};
				\pic[draw,thin,angle radius=2mm] {right angle = A--H--B};
				\pic[draw,thin,angle radius=2mm] {right angle = C--B--A};
				\foreach \i/\g in {S/90,A/-135,B/-90,C/-45,H/45}{\draw[fill=black](\i) circle (1pt) ($(\i)+(\g:3mm)$) node[scale=1]{$\i$};}
			\end{tikzpicture}
		\end{center}
		Trong mặt phẳng $(SAB)$, kẻ $AH\perp SB$.\\
		Ta có $\heva{&BC\perp AB\\&BC\perp SA}$ suy ra $BC\perp(SAB)$.\\
		Mà $AH\subset(SAB)$ nên $BC\perp AH$.\\
		Từ $(1)$ và $(2)$ suy ra $AH\perp(SBC)$ hay $\mathrm{d}(A, (SBC))=AH$.\\
		Xét $\triangle SAB$ vuông tại $A$ có 
		$$\dfrac{1}{AH^2}=\dfrac{1}{AS^2}+\dfrac{1}{AB^2}=\dfrac{1}{4a^2}+\dfrac{1}{a^2}=\dfrac{5}{4a^2} \Rightarrow AH=\dfrac{2a \sqrt{5}}{5}.$$
		Vậy $\mathrm{d}(A, (SBC))=AH=\dfrac{2a \sqrt{5}}{5}$.
	}
\end{ex}

\begin{ex}%[1H4B5-4]
	Cho hình chóp $S.ABCD$ có đáy là hình chữ nhật với $AB=a$, $AD=2a$. Biết $SA=2a$ và $SA$ vuông góc với mặt phẳng đáy. Khoảng cách giữa hai đường thẳng $BD$ và $SC$ bằng
	\choice
	{$\dfrac{a \sqrt{30}}{6}$}
	{\True $\dfrac{2\sqrt{21} a}{21}$}
	{$\dfrac{4\sqrt{21} a}{21}$}
	{$\dfrac{a \sqrt{30}}{12}$}
	\loigiai{
		\begin{center}
			\begin{tikzpicture}[scale=1, font=\footnotesize,line join=round, line cap=round, >=stealth]
				\def\a{4} % Chiều dài mặt đáy
				\def\b{2} % Chiều rộng mặt đáy
				\def\g{35} % Góc xiên
				\def\h{3.5} % Chiều cao hình chóp
				\path  (0,0) coordinate (B)
				(0:\a) coordinate (C)
				(\g:\b) coordinate (A)
				(A)+(0:\a) coordinate (D)
				($(B)!0.5!(D)$) coordinate (O)
				(A)+(90:\h) coordinate (S)
				($(S)!0.5!(A)$) coordinate (M)
				;
				\foreach \i in {B,C,D}{\draw (S)--(\i);}
				\draw (B)--(C)--(D);
				\draw[dashed] (S)--(A)--(B)--(D)--(A)--(C) (B)--(M)--(D) (O)--(M);
				\pic[draw,thin,angle radius=2mm] {right angle = C--A--S};
				\foreach \i/\g in {S/90,A/135,B/-135,C/-45,D/45,O/-90,M/170}{\draw[fill=black](\i) circle (1pt) ($(\i)+(\g:3mm)$) node[scale=1]{$\i$};}
			\end{tikzpicture}
		\end{center}
		Gọi $O$ là tâm hình chữ nhật và $M$ là trung điểm $SA$, ta có $SC\parallel(BMD)$.\\
		Do đó $\mathrm{d}(SC, BD)=\mathrm{d}(SC, (BMD))=\mathrm{d}(S, (BMD))=\mathrm{d}(A, (BMD))=h$.\\
		Ta có $AM$, $AB$, $AD$ đôi một vuông góc nên
		$$
		\dfrac{1}{h^2}=\dfrac{1}{AM^2}+\dfrac{1}{AB^2}+\dfrac{1}{AD^2}=\dfrac{4}{a^2}+\dfrac{1}{a^2}+\dfrac{1}{4a^2}.
		$$
		Vậy $\mathrm{d}(SC, BD)=\dfrac{2a \sqrt{21}}{21}$.
	}
\end{ex}
\Closesolutionfile{ans}
\indapan{6}{ans/ans-De1}
\cauds
\Opensolutionfile{ans}[ans/ans-De1-DS]
\begin{ex}%[1H4K7-3]
	Cho hình chóp $S.ABCD$ có đáy $ABCD$ là hình chữ nhật với $AB=a \sqrt{2}$, $\linebreak AC=a \sqrt{3}$. Cạnh bên $SA=2a$ và vuông góc với mặt đáy $(ABCD)$. Khi đó
	\choiceTF
	{\True $AD\parallel (SBC)$}
	{Khoảng cách từ $D$ đến mặt phẳng $(SBC)$ bằng $\dfrac{a\sqrt{3}}{3}$}
	{\True Khoảng cách giữa hai đường thẳng ${SD,AB}$ bằng $\dfrac{2a\sqrt{5}}{5}$}
	{Thể tích khối chóp $S.ABCD$ bằng $\dfrac{\sqrt{2}a^3}{3}$}
	\loigiai{
		\begin{center}
			\begin{tikzpicture}[scale=1, font=\footnotesize,line join=round, line cap=round, >=stealth]
				\def\a{4} % Chiều dài mặt đáy
				\def\b{2} % Chiều rộng mặt đáy
				\def\g{35} % Góc xiên
				\def\h{3.5} % Chiều cao hình chóp
				\path  (0,0) coordinate (B)
				(0:\a) coordinate (C)
				(\g:\b) coordinate (A)
				(A)+(0:\a) coordinate (D)
				($(S)!(A)!(B)$) coordinate (H)
				($(S)!(A)!(D)$) coordinate (K)
				(A)+(90:\h) coordinate (S)
				;
				\foreach \i in {B,C,D}{\draw (S)--(\i);}
				\draw (B)--(C)--(D);
				\draw[dashed] (S)--(A)--(B) (D)--(A) (H)--(A)--(K);
				\pic[draw,thin,angle radius=2mm] {right angle = B--H--A};
				\pic[draw,thin,angle radius=2mm] {right angle = A--K--D};
				\foreach \i/\g in {S/90,A/135,B/-135,C/-45,D/45,H/160,K/45}{\draw[fill=black](\i) circle (1pt) ($(\i)+(\g:3mm)$) node[scale=1]{$\i$};}
			\end{tikzpicture}
		\end{center}
		\begin{itemchoice}
			\itemch Đúng. Ta có $AD\parallel BC\Rightarrow AD\parallel (SBC)$.
			\itemch Sai. Vì $AD\parallel (SBC) \Rightarrow \mathrm{d}(D, (SBC))=\mathrm{d}(A, (SBC))$.\\
			Trong mặt phẳng $(SAB)$, kẻ ${AH\perp SB}$ tại $H$. \hfill(1)\\
			Ta có $\heva{&BC\perp AB\\ &BC\perp SA} \Rightarrow BC\perp(SAB) \Rightarrow AH\perp BC$. \hfill(2)\\
			Từ $(1)$ và $(2)$ suy ra $AH\perp(SBC)$ hay $\mathrm{d}(A, (SBC))=AH$.\\
			Tam giác $SAB$ vuông tại $A$ có đường cao $AH$ nên
			$$\dfrac{1}{AH^2}=\dfrac{1}{SA^2}+\dfrac{1}{AB^2} \Rightarrow AH=\dfrac{SA\cdot AB}{\sqrt{SA^2+AB^2}}=\dfrac{2a \cdot a \sqrt{2}}{\sqrt{4a^2+2a^2}}=\dfrac{2a \sqrt{3}}{3}.$$
			Vậy $\mathrm{d}(D,(SBC))=\mathrm{d}(A,(SBC))=AH=\dfrac{2a\sqrt{3}}{3}$.
			\itemch Đúng. Trong mặt phẳng $(SAD)$, kẻ $AK\perp SD$ tại $K$. \hfill(3)\\
			Ta có $\heva{&AB\perp SA\\ &AB\perp AD} \Rightarrow AB\perp(SAD) \Rightarrow AB\perp AK$.	\hfill(4)\\
			Từ $(3)$ và $(4)$ suy ra $AK$ là đường vuông góc chung của hai đường thẳng chéo nhau ${AB, SD}$.\\
			Tam giác $ACD$ vuông tại $D$ nên $AD=\sqrt{AC^2-CD^2}=\sqrt{3a^2-2a^2}=a$.\\
			Tam giác $SAD$ vuông tại $A$ có đường cao $AK$ nên
			$$\dfrac{1}{AK^2}=\dfrac{1}{SA^2}+\dfrac{1}{AD^2} \Rightarrow AK=\dfrac{SA\cdot AD}{\sqrt{SA^2+AD^2}}=\dfrac{2a\cdot a}{\sqrt{4a^2+a^2}}=\dfrac{2a \sqrt{5}}{5}.$$
			Vậy $\mathrm{d}(AB, SD)=AK=\dfrac{2a\sqrt{5}}{5}$.
			\itemch Sai. Diện tích đáy hình chóp là $S_{ABCD}=a \cdot a \sqrt{2}=a^2\sqrt{2}$.\\
			Thể tích khối chóp cần tìm là
			$V_{S.ABCD}=\dfrac{1}{3}SA\cdot S_{ABCD}=\dfrac{1}{3}\cdot 2a\cdot a^2\sqrt{2}=\dfrac{2\sqrt{2}a^3}{3}$ (đvtt).
		\end{itemchoice}
	}
\end{ex}

\begin{ex}%[1H4K5-4]
	Cho hình chóp $S.ABCD$ có đáy $ABCD$ là hình chữ nhật với $AB=2a$, $AD=a$. Hình chiếu của $S$ lên mặt phẳng $(ABCD)$ là trung điểm $H$ của $AB$ và $\widehat{SCH}=45^{\circ}$. Khi đó
	\choiceTF
	{\True $BC\perp (SAB)$}
	{\True $\mathrm{d}(H,(SBC))=\dfrac{a\sqrt{6}}{3}$}
	{\True Gọi $K$ là trung điểm $CD$ khi đó $CD\perp (SHK)$}
	{$\mathrm{d}(H, (SCD))=\dfrac{a\sqrt{6}}{2}$}
	\loigiai{
		\begin{center}
			\begin{tikzpicture}[scale=1, font=\footnotesize,line join=round, line cap=round, >=stealth]
				\def\a{4} % Chiều dài mặt đáy
				\def\b{2.5} % Chiều rộng mặt đáy
				\def\g{55} % Góc xiên
				\def\h{3.5} % Chiều cao hình chóp
				\path  (0,0) coordinate (B)
				(0:\a) coordinate (C)
				(\g:\b) coordinate (A)
				(A)+(0:\a) coordinate (D)
				($(A)!0.5!(B)$) coordinate (H)	
				($(C)!0.5!(D)$) coordinate (K)	
				(H)+(90:\h) coordinate (S)
				($(B)!0.3!(S)$) coordinate (E)
				($(S)!(H)!(K)$) coordinate (I)
				;
				\foreach \i in {B,C,D}{\draw (S)--(\i);}
				\draw (B)--(C)--(D) (S)--(K);
				\draw[dashed] (S)--(A)--(B) (D)--(A) (S)--(H)--(E) (C)--(H)--(I) (H)--(K);
				\pic[draw,thin,angle radius=2mm] {right angle = A--H--S};
				\pic[draw,thin,angle radius=2mm] {right angle = H--K--C};
				\pic[draw,thin,angle radius=2mm] {right angle = H--I--K};
				\pic[draw,thin,angle radius=2mm] {right angle = B--E--H};
				\pic[draw, angle radius=15pt,fill=gray!20]{angle= S--C--H};
				\foreach \i/\g in {S/90,A/135,B/-135,C/-45,D/45,E/160,K/-45,I/45,H/-80}{\draw[fill=black](\i) circle (1pt) ($(\i)+(\g:3mm)$) node[scale=1]{$\i$};}
			\end{tikzpicture}
		\end{center}
		\begin{itemchoice}
			\itemch Đúng. Kẻ đường cao $HE$ trong tam giác $SBH$. \hfill(1)\\
			Ta có $\heva{&BC\perp AB\\ &BC\perp SH\ (\text{do } SH\perp(ABCD))} \Rightarrow BC\perp(SAB)$.
			\itemch Đúng. Ta có $BC\perp(SAB) \Rightarrow BC\perp HE$. \hfill(2)\\
			Từ $(1)$ và $(2)$ suy ra $HE\perp(SBC)$ hay $\mathrm{d}(H,(SBC))=HE$.\\
			Tam giác $BCH$ vuông tại $B$ có
			$$CH=\sqrt{BC^2+BH^2}=\sqrt{a^2+a^2}=a \sqrt{2}.$$
			Tam giác $SCH$ vuông tại $H$ có
			$\tan \widehat{SCH}=\dfrac{SH}{CH} \Rightarrow SH=a \sqrt{2}$.\\
			Tam giác $SBH$ vuông tại $H$ có đường cao $HE$ nên
			\begin{eqnarray*}
				\dfrac{1}{HE^2} &=&\dfrac{1}{SH^2}+\dfrac{1}{BH^2} \\
				\Rightarrow HE&=&\dfrac{SH\cdot BH}{\sqrt{SH^2+BH^2}} \\
				&=&\dfrac{a \sqrt{2} \cdot a}{\sqrt{2 a^2+a^2}}=\dfrac{a \sqrt{6}}{3}.
			\end{eqnarray*}
			Vậy $\mathrm{d}(H, (SBC))=HE=\dfrac{a \sqrt{6}}{3}$.
			\itemch Đúng. Gọi $K$ là trung điểm $CD$ thì $HK$ là đường trung bình của hình chữ nhật $ABCD$ nên ${HK\parallel BC\parallel AD\Rightarrow HK\perp CD}$.\\
			Ta có ${\heva{&CD\perp HK\\ &CD\perp SH} \Rightarrow CD\perp(SHK)}$.
			\itemch Sai. Kẻ đường cao $HI$ của tam giác $SHK$.\\
			Ta có $\heva{&HI\perp SK\\ &HI\perp CD\ (\text{do } CD\perp(SHK), HI\subset(SHK))} \Rightarrow HI\perp(SCD)$.\\
			Tam giác $SHK$ vuông tại $H$ có đường cao $HI$ nên
			$$\dfrac{1}{HI^2}=\dfrac{1}{SH^2}+\dfrac{1}{HK^2} \Rightarrow HI=\dfrac{SH\cdot HK}{\sqrt{SH^2+HK^2}}=\dfrac{a \sqrt{2} \cdot a}{\sqrt{2a^2+a^2}}=\dfrac{a \sqrt{6}}{3}.$$
			Vậy $\mathrm{d}(H, (SCD))=HI=\dfrac{a \sqrt{6}}{3}$.
		\end{itemchoice}
	}
\end{ex}

\begin{ex}%[1H4B5-4]
	Cho hình chóp đều $S.ABC$ có cạnh đáy bằng $a$, gọi $O$ là tâm của đáy và $SO=\dfrac{a \sqrt{3}}{3}$. Khi đó
	\choiceTF
	{$AO=\dfrac{a\sqrt{3}}{2}$}
	{\True $\mathrm{d}(O, SA)=\dfrac{a\sqrt{6}}{6}$}
	{\True Kẻ đường cao $AI$ của tam giác $ABC$, khi đó $OI=\dfrac{a\sqrt{3}}{6}$}
	{$\mathrm{d}(O,(SBC))=\dfrac{a\sqrt{15}}{12}$}
	\loigiai{
		\begin{center}
			\begin{tikzpicture}[scale=1, font=\footnotesize,line join=round, line cap=round, >=stealth]
				\def\b{4.5}  % Chiều dài AC
				\def\c{3.5}  % Chiều dài AB
				\def\gA{-30} % Góc A
				\def\h{3.5}  % Chiều cao hình chóp
				\path  (0,0) coordinate (A)
				(0:\b) coordinate (C)
				(\gA:\c) coordinate (B)
				($(B)!0.5!(C)$) coordinate (I)				% Lấy trung điểm
				($(A)!2/3!(I) $)			coordinate(O) 	% ($(O)!k!(A) $) Vị tự điểm A gốc O k lần OA
				(O)+(90:\h) coordinate (S)
				($(S)!(O)!(A)$) coordinate (H)
				($(S)!(O)!(I)$) coordinate (K)
				;
				\foreach \i in {A,B,C}{\draw (S)--(\i);}
				\draw (A)--(B)--(C) (S)--(I);
				\draw[dashed] (A)--(C) (A)--(I) (H)--(O)--(K) (S)--(O);
				\pic[draw,thin,angle radius=2mm] {right angle = O--H--S};
				\pic[draw,thin,angle radius=2mm] {right angle = S--O--A};
				\pic[draw,thin,angle radius=2mm] {right angle = S--K--O};
				\pic[draw,thin,angle radius=2mm] {right angle = A--I--B};
				\foreach \i/\g in {S/90,A/-135,B/-90,C/-45,O/-90,H/150,K/0,I/0}{\draw[fill=black](\i) circle (1pt) ($(\i)+(\g:3mm)$) node[scale=1]{$\i$};}
			\end{tikzpicture}
		\end{center}
		
		\begin{itemchoice}
			\itemch Sai. Kẻ đường cao $AI$ của tam giác $ABC$, ta có $O$ thuộc $AI$.\\
			Trong mặt phẳng $(SAI)$, dựng $OH\perp SA$ tại $H\Rightarrow \mathrm{d}(O, SA)=OH$.\\
			Tam giác $ABC$ đều cạnh $a$ nên $AO=\dfrac{2}{3} AI=\dfrac{2}{3} \cdot \dfrac{a \sqrt{3}}{2}=\dfrac{a \sqrt{3}}{3}=SO$.
			\itemch Đúng. Tam giác $SAO$ vuông cân tại $O$ nên $OH=\dfrac{SA}{2}=\dfrac{\dfrac{a \sqrt{3}}{3} \cdot \sqrt{2}}{2}=\dfrac{a \sqrt{6}}{6}$.\\
			Vậy $\mathrm{d}(O, SA)=OH=\dfrac{a\sqrt{6}}{6}$.
			\itemch Đúng. Ta xét khoảng cách từ $O$ đến mặt bên $(SBC)$.\\
			Kẻ đường cao $OK$ của tam giác $SOI\Rightarrow OK\perp SI$. \hfill(1)\\
			Ta có $\heva{&BC\perp SO\ (\text{do } SO\perp(ABC) \\ &BC\perp AI} \Rightarrow BC\perp(SAI) \Rightarrow BC\perp OK$. \hfill(2)\\
			Từ $(1)$ và $(2)$ suy ra $OK\perp(SBC)$ hay $OK=\mathrm{d}(O, (SBC))$.\\
			Ta có ${OI=\dfrac{AI}{3}=\dfrac{\dfrac{a \sqrt{3}}{2}}{3}=\dfrac{a \sqrt{3}}{6}}$.
			\itemch Sai. Tam giác $SOI$ vuông tại $O$ có đường cao $OK$ nên
			$$\dfrac{1}{OK^2}=\dfrac{1}{SO^2}+\dfrac{1}{OI^2} \Rightarrow OK=\dfrac{SO\cdot OI}{\sqrt{SO^2+OI^2}}=\dfrac{\dfrac{a \sqrt{3}}{3} \cdot \dfrac{a \sqrt{3}}{6}}{\sqrt{\dfrac{3a^2}{9}+\dfrac{3a^2}{36}}}=\dfrac{a \sqrt{15}}{15}.$$
			Vậy ${\mathrm{d}(O, (SBC.)=OK=\dfrac{a \sqrt{15}}{15}}$.
		\end{itemchoice}
	}
\end{ex}
\begin{ex}%[1H8H5-3]
	Cho hình chóp $S.ABCD$ có $SA\perp(ABCD)$, $SA=a\sqrt{3}$, $ABCD$ là hình vuông cạnh bằng $a$. Khi đó:
	\choiceTF
	{$\mathrm{d}(A,(SBC))=\dfrac{\sqrt{3}}{3}a$}
	{\True $AD\parallel(SBC)$}
	{\True $\mathrm{d}(D,(SBC))=\dfrac{\sqrt{3}}{2}a$}
	{\True Gọi $M$ là trung điểm của $SA$. Khi đó: $\mathrm{d}(M,(SBC))=\dfrac{\sqrt{3}}{4}a$}
	\loigiai{
		\begin{center}
			\begin{tikzpicture}
				\def\a{4}
				\def\h{4}
				\path 	(0:0) coordinate (A)
				++(0:\a) coordinate (D)
				++(-130:\a/2) coordinate (C)
				($(A)+(C)-(D)$) coordinate (B)
				($(A)+(90:\h)$) coordinate (S)
				($(S)!0.5!(A)$) coordinate (M)
				($(S)!0.5!(B)$) coordinate (H)
				(intersection of A--C and B--D) coordinate (O);%giao điểm O
				\draw[dashed,thick] 	(B)--(A)--(D)	(H)--(A)--(S) (C)--(A) (B)--(D);
				\draw[thick] 			(B)-- (C)--(D)
				(B)--(S)	(C)--(S)	(D)--(S);
				\foreach \x/\g in {A/180,B/-135,C/-45,D/45,S/90,M/0,H/135,O/90}
				\fill[black] 	(\x) circle (1.5pt)
				($(\g:3mm)+(\x)$) node {$\x$};	
				\draw pic[draw,angle radius=3mm]{right angle=D--A--S};%Theo chiều dương	
				\draw pic[draw,angle radius=3mm]{right angle=A--H--S};%Theo chiều dương	
			\end{tikzpicture}
		\end{center}
		Kẻ $SH\perp SB$ tại $H$.\\
		Ta có  $\heva{&BC\perp SA\\&BC\perp AB}\Rightarrow BC\perp (SAB) \Rightarrow BC\perp AH$.\\
		Ta lại có $AH\perp SB \Rightarrow AH\perp (SBC)\Rightarrow \mathrm{d}(A,(SBC))=AH$.\\
		Ta có  $\dfrac{1}{AH^2}=\dfrac{1}{SA^2}+\dfrac{1}{AB^2}=\dfrac{1}{3a^2}+\dfrac{1}{a^2}\Rightarrow AH=\dfrac{a\sqrt{3}}{2}$.\\
		Vậy $\mathrm{d}(A,(SBC))=\dfrac{\sqrt {3}}{2}a$.\\
		Ta có $AD\parallel(SBC) \Rightarrow \mathrm{d}(D,(SBC))=\dfrac{\sqrt{3}}{2}a$.\\
		Ta có $MA$ cắt $(SBC)$ tại $S$. \\
		Suy ra $ \dfrac{\mathrm{d}(M,(SBC))}{\mathrm{d}(A,(SBC))}=\dfrac{MS}{AS}=\dfrac{1}{2}\Rightarrow \mathrm{d}(M,(SBC))=\dfrac{1}{2}\mathrm{d}(A,(SBC))=\dfrac{1}{2} \cdot\dfrac{\sqrt {3}}{2}a=\dfrac{\sqrt{3}}{4}a$. 
		\begin{itemchoice}
			\itemch Sai. Vì $\mathrm{d}(A,(SBC))=\dfrac{\sqrt {3}}{2}a$.
			\itemch  Đúng. Vì $AD\parallel(SBC)$.
			\itemch Đúng. Vì $\mathrm{d}(D,(SBC))=\dfrac{\sqrt{3}}{2}a$.
			\itemch Đúng. Vì $\mathrm{d}(M,(SBC))=\dfrac{\sqrt{3}}{4}a$
		\end{itemchoice}
	}
\end{ex}
\Closesolutionfile{ans}
\indapan{3}{ans/ans-De1-DS}
\Opensolutionfile{ans}[ans/ans-De1-KQ]
\caukq
\begin{ex}%[1H8V5-3]
	Cho hình chóp $S.ABCD$ có đáy là hình vuông cạnh $a$, tam giác $SAB$  đều và nằm trong mặt phẳng vuông góc với đáy. Với $a=\sqrt{21}$, khoảng cách từ $A$ đến $(SDC)$ là
	\shortans{$3$}
	\loigiai{
		\begin{center}
			\begin{tikzpicture}
				\def\a{4}
				\def\h{4.5}
				\path 	(0:0) coordinate (A)
				++(0:\a) coordinate (D)
				++(-130:\a/2) coordinate (C)
				($(A)+(C)-(D)$) coordinate (B)
				($(A)!0.5!(B)$) coordinate (H)
				($(C)!0.5!(D)$) coordinate (K)
				($(H)+(90:\h)$) coordinate (S)
				($(S)!0.6!(K)$) coordinate (I)
				(intersection of A--C and B--D) coordinate (O);%giao điểm O
				\draw[dashed,thick] 	(B)--(A)--(D)	(A)--(S) (H)--(I)
				(S)--(H)--(K)	;
				\draw[thick] 			(B)--(C)--(D)
				(B)--(S)--(K)	(C)--(S)	(D)--(S);
				\draw pic[draw,angle radius=2mm]{right angle=H--I--S};%Theo chiều dương
				\foreach \x/\g in {A/135,B/-135,C/-45,D/45,S/90,H/135,K/-90,I/90}
				\fill[black] 	(\x) circle (1.5pt)
				($(\g:4mm)+(\x)$) node {$\x$};
			\end{tikzpicture}
		\end{center}
		Gọi $H$ là trung điểm $AB$, suy ra $SH\perp AB$  (do tam giác $SAB$ đều).\\
		Mặt khác $(SAB)\perp (ABCD)$  nên $SH\perp (ABCD)$.\\
		Ta có $AH\parallel CD$ nên $AH\parallel (SCD)$.\\
		Kẻ $HK\perp CD$ và $HI\perp SK$, suy ra $\mathrm{d}(H,(SCD))=HI$.\\
		Nên $\mathrm{d}(A,(SCD))=\mathrm{d}(H,(SCD))=HI$.\\
		Ta có $HK=a$, đường cao hình chóp là $SH=\dfrac{a\sqrt{3}}{2}$.\\
		Xét tam giác $SHK$ vuông tại $H$ có $HI$ là đường cao, nên ta có\\ $\dfrac{1}{HI^2}=\dfrac{1}{HK^2}+\dfrac{1}{HS^2}\Rightarrow \dfrac{1}{HI^2}=\dfrac{1}{a^2}+\dfrac{4}{3a^2}\Rightarrow HI=\dfrac{a\sqrt{21}}{7}=3$ với $a=\sqrt{21}$.
	}
\end{ex}

\begin{ex}%[1H8H5-2]
	Cho tứ diện $S.ABC$  trong đó $SA,SB,SC$ vuông góc với nhau từng đôi một và $SA=3a$, $SB=a$, $SC=2a$. Với $a=\sqrt{5}$, tính khoảng cách từ $A$ đến đường thẳng $BC$.
	\shortans{$7$}
	\loigiai{
		
		\begin{center}
			\begin{tikzpicture}
				\def\a{4} %Khai báo cạnh
				\def\h{4}
				\path 	(0:0) coordinate (S)
				++(0:\a) coordinate (B)
				++(-150:4*\a/5) coordinate (C)
				($(S)+(90:\h)$) coordinate (A)
				($(S)!0.5!(B)$) coordinate (M)
				($(B)!0.5!(C)$) coordinate (H);
				\draw[thick] 	(S)--(C)--(B)
				(S)--(A)	(B)--(A)--(H)	(C)--(A);
				\draw[dashed,thick] 	(H)--(S)--(B);
				\foreach \x /\goc in {A/180,B/0,C/-135,S/-135,H/-20}
				\fill[black] (\x) circle (1.5pt)
				($(\x)+(\goc:3mm)$) node {$\x$};
				\draw pic[draw,angle radius=2mm]{right angle=B--S--A};%Theo chiều dương
				\draw pic[draw,angle radius=2mm]{right angle=A--H--B};%Theo chiều dương
			\end{tikzpicture}
		\end{center}
		Kẻ $AH\perp BC$ tại $H$ Suy ra $\mathrm{d}(A,BC)=AH$.\\
		Ta có $\heva{&BC\perp SA\\&BC\perp AH} \Rightarrow BC\perp(SAH)\Rightarrow BC\perp SH$.\\
		Ta có $SH =\dfrac{1}{\sqrt{\dfrac{1}{SC^2}+\dfrac{1}{SB^2}}}=\dfrac{1}{\sqrt{\dfrac{1}{{(2a)}^2}+\dfrac{1}{a^2}}}=\dfrac{2\sqrt {5}}{5}a$.\\
		Ta có $AH=\sqrt{SA^2+SH^2}=\sqrt{\left(3a\right)^2+\left(\dfrac{2\sqrt{5}}{5}a\right)^2}=\dfrac{7\sqrt{5}}{5}a$.\\
		Vậy $\mathrm{d}(A,BC)=\dfrac{7\sqrt{5}}{5}a=7$ với $a=\sqrt{5}$.
	}
\end{ex}

\begin{ex}%[1H8V5-3]
	Cho hình chóp đều $S.ABC$ có đáy cạnh $a$ và cạnh bên $2a$. Với $a=10$, tính khoảng cách từ $A$ đến mặt phẳng $(SBC)$ (kết quả làm tròn hai chữ số thập phân).
	\shortans{$8{,}56$}
	\loigiai{
		\begin{center}
			\begin{tikzpicture}
				\def\a{4}
				\path 	(0:0) coordinate (B)
				++(0:\a) coordinate (A)
				++(-120:\a/2) coordinate (C)
				($(B)+(70:\a)$) coordinate (S)
				($(B)!0.5!(C)$) coordinate (I)
				($(S)!0.8!(I)$) coordinate (H)
				($(A)!0.7!(I)$) coordinate (O);
				\draw[dashed,thick] 	(B)--(A)--(I) (O)--(H);
				\draw[thick] 			(B)--(S)--(A) (I)--(S)--(C)
				(B)--(C)--(A);
				\foreach \x/\g in {S/90,B/180,C/-45,A/0,I/-90,H/135,O/-45}
				\fill[black] 	(\x) circle (1pt)
				($(\g:3mm)+(\x)$) node {$\x$};
				\draw pic[draw,angle radius=2mm]{right angle=O--H--S};%Theo chiều dương
				\draw pic[draw,angle radius=2mm]{right angle=A--I--C};%Theo chiều dương
				%Hình chóp S.ABC có SA vuông góc đáy	
			\end{tikzpicture}
		\end{center}
		
		Kẻ $OH\perp SI$ tại $H$. Ta có $\heva{&BC\perp OI\\&BC\perp SO}\Rightarrow BC\perp (SOI)\Rightarrow BC\perp OH$.\\
		Ta lại có $OH\perp SI \Rightarrow OH\perp (SBC)\Rightarrow \mathrm{d}(O,(SBC))=OH$.\\
		Ta có $OI=\dfrac{1}{3}\cdot\dfrac{a\sqrt{3}}{2}=\dfrac{\sqrt{3}}{6}a$; $SO=\sqrt{SB^2-OB^2}=\sqrt{(2a)^2-\left(\dfrac{2}{3}\cdot\dfrac{a\sqrt{3}}{2}\right)^2}=\dfrac{\sqrt{33}}{3}a$.\\
		
		$OH=\dfrac{1}{\sqrt{\dfrac{1}{SO^2}+\dfrac{1}{OI^2}}}=\dfrac{1}{\sqrt{\dfrac{1}{\left(\dfrac{\sqrt{33}}{3}a\right)^2}+\dfrac{1}{\left(\dfrac{\sqrt{3}}{6}a\right)^2}}}=\dfrac{\sqrt{165}}{45}a$.\\
		Vậy $\mathrm{d}(O,(SBC))=\dfrac{\sqrt{165}}{45}a$.\\
		Ta có $AO$ cắt $(SBC)$ tại $I$.\\
		$\Rightarrow \dfrac{\mathrm{d}(A,(SBC))}{\mathrm{d}(O,(SBC)}=\dfrac{AI}{OI}=3\Rightarrow \mathrm{d}(A,(SBC))=3\mathrm{d}(O,(SBC))=\dfrac{\sqrt {165}}{15}a \approx 8{,}56$ với $a=10$.	
	}
\end{ex}

\begin{ex}%[1H8V5-3]
	Cho hình chóp $S.ABC$ có đáy là tam giác đều cạnh $a$, $SA\perp (ABC)$ và $SB=2a$. Gọi $G$ là trọng tâm tam giác $ABC$. Với $a=\sqrt{15}$, tính khoảng cách từ $G$ đến mặt phẳng $(SBC)$.
	\shortans{$1$}
	\loigiai{
		\begin{center}
			\begin{tikzpicture}
				\def\a{4} %Khai báo cạnh
				\def\h{4}
				\path 	(0:0) coordinate (A)
				++(0:\a) coordinate (B)
				++(-150:4*\a/5) coordinate (C)
				($(A)+(90:\h)$) coordinate (S)
				($(A)!0.5!(B)$) coordinate (M)
				($(C)!0.5!(B)$) coordinate (I)
				($(S)!0.3!(I)$) coordinate (H)
				($(A)!2/3!(I)$) coordinate (G);
				\draw[thick] 	(A)--(C)--(B)
				(A)--(S)	(B)--(S)--(I)	(C)--(S);
				\draw[dashed,thick] 	(I)--(A)--(B) (H)--(A);
				\foreach \x /\goc in {A/180,B/0,C/-135,S/90,I/-20,G/-90,H/0}
				\fill[black] (\x) circle (1.5pt)
				($(\x)+(\goc:3mm)$) node {$\x$};
				\draw pic[draw,angle radius=2mm]{right angle=A--I--C};%Theo chiều dương
				\draw pic[draw,angle radius=2mm]{right angle=A--H--I};%Theo chiều dương
			\end{tikzpicture}
		\end{center}
		Kẻ $AI\perp BC$, kẻ $AH\perp SI$ tại $H$.\\
		Ta có $\heva{&BC\perp SA\\&BC\perp AI}\Rightarrow BC\perp (SAI)\Rightarrow BC\perp AH$.\\
		Ta lại có $AH\perp SI \Rightarrow AH\perp (SBC) \Rightarrow \mathrm{d}(A,(SBC))=AH$.\\
		Ta có $SA=\sqrt{SB^2-BA^a2}=\sqrt{(2a)^2-a^2}=\sqrt{3}a$.\\
		Ta có $\dfrac{1}{AH^2}=\dfrac{1}{SA^2}+\dfrac{1}{AI^2}=\dfrac{1}{3a^2}+\dfrac{4}{3a^2}\Rightarrow AH=\dfrac{\sqrt{15}}{5}a$.\\
		Vậy $\mathrm{d}(A,(SBC))=\dfrac{\sqrt{15}}{5}a$.\\
		Ta có $GA$ cắt $BC$ tại $I$.\\
		$\Rightarrow\dfrac{\mathrm{d}(G,(SBC))}{{\mathrm{d}(A,(SBC))}} =\dfrac{GI}{AI} =\dfrac{1}{3}\Rightarrow \mathrm{d}(G,(SBC))=\dfrac{1}{3}\mathrm{d}(A,(SBC))=\dfrac{\sqrt{15}}{15}a=1$ với $a=\sqrt{15}$. 
	}
\end{ex}

\begin{ex}%[1H8H5-3]
	Cho hình hộp chữ nhật $ABCD.A'B'C'D'$ có ba kích thước $AB=7$, $AD=14$, $AA'=21$. Khoảnh cách từ $A$ đến mặt phẳng $(A'BD)$ bằng bao nhiêu?
	\shortans{$6$}
	\loigiai{
		\begin{center}
			\begin{tikzpicture}
				\def\bc{4} % cạnh BC
				\def\ba{2} % cạnh BA
				\def\gocB{35} % góc B của đáy
				\coordinate[label=below left:$B$] (B) at (0,0);
				\coordinate[label=above left:$A$] (A) at (\gocB:\ba);
				\coordinate[label=below:$C$] (C) at (\bc,0);
				\coordinate[label=right:$D$] (D) at ($(C)-(B)+(A)$);
				\coordinate[label=above left:$A'$] (A') at ($(A)+(90:\bc)$);
				\coordinate[label=left:$B'$] (B') at ($(B)-(A)+(A')$);
				\coordinate[label=below right:$C'$] (C') at ($(C)-(A)+(A')$);
				\coordinate[label=right:$D'$] (D') at ($(D)-(A)+(A')$);
				\coordinate[label=below right:$M$] (M) at ($(B)!0.5!(D)$);
				\coordinate[label=right:$K$] (K) at ($(A')!0.7!(M)$);
				\draw (B')--(B)--(C)--(D)--(D')--(A')--(B')--(C')--(D') (C)--(C');
				\draw[dashed] (A')--(A)--(D) (A)--(B)--(D) (K)--(A)--(M)--(A');
				\draw pic[draw,angle radius=2mm]{right angle=A--K--M};%Theo chiều dương
				\foreach \diem in {A,B,C,D,A',B',C',D',M,K}	\fill (\diem)circle(1.5pt);
			\end{tikzpicture}
		\end{center}
		Đặt $AB=a$, $AD=2a$, $AA'=3a$.\\
		Trong tam giác $ABD$ kẻ $AM\perp BD$ suy ra $\dfrac{1}{AM^2}=\dfrac{1}{AB^2}+\dfrac{1}{AD^2}$.\\
		Trong tam giác $A'AM$ kẻ $AK\perp A'M$ $\Rightarrow \dfrac{1}{AK^2}=\dfrac{1}{A'A^2}+\dfrac{1}{AM^2}=\dfrac{1}{a^2}+\dfrac{1}{4a^2}+\dfrac{1}{9a^2}$.\\
		$\Rightarrow AK=\dfrac{6}{7}a$.\\
		Khi đó $AK\perp \left(A'BD\right)$ hay $\mathrm{d}\left( A,\left(A'BD\right)\right)=AK=\dfrac{6}{7}a=6$ với $a=7$.
	}
\end{ex}
\begin{ex}%[1H8H5-3]
	Cho khối lăng trụ đều $ABC.A'B'C'$ có cạnh đáy bằng $\sqrt{5}$, $A'A=2\sqrt{5}$. Khoảng cách từ điểm $A'$ đến mặt phẳng $(AB'C')$ bằng  
	\shortans{$2$}
	\loigiai{
		
		\begin{center}
			\begin{tikzpicture}
				\def\a{4}
				\def\h{4.5}
				\path 	(0:0) coordinate (A)
				++(0:\a) coordinate (C)
				++(-150:3*\a/4) coordinate (B)
				($(A)+(90:\h)$) coordinate (A')
				($(B)+(90:\h)$) coordinate (B')
				($(C)!0.5!(B)$) coordinate (M)
				($(A')!0.4!(M)$) coordinate (H)
				($(C)+(90:\h)$) coordinate (C');
				\draw[dashed,thick] 	(H)--(A)--(C)--(A')--(M)--(A);
				\draw[thick]	(C)--(C') 	(B)--(B')	(A)--(A') (A)--(B)--(C) (A')--(B) (A')--(B')--(C')--cycle;
				\draw pic[draw,angle radius=2mm]{right angle=A--H--M};%Theo chiều dương
				\foreach \x/\g in {A/180,B/-45,C/0,A'/180,B'/90,C'/0,M/-20}
				\fill[black] 	(\x) circle (1pt)
				($(\g:4mm)+(\x)$) node {$\x$};	
			\end{tikzpicture}
		\end{center}
		Gọi cạnh đáy là $a \Rightarrow AA'=2a$.\\
		Gọi $M$ là trung điểm của $B'C'$.\\
		Ta có  $\heva{&AA'\perp B'C'\\&A'M\perp B'C'}\Rightarrow B'C'\perp \left(AA'M\right)\Rightarrow\left(AB'C'\right) \perp \left(AA'M\right)$
		theo giao tuyến $AM$.\\
		Kẻ $A'H\perp AM$ trong mặt phẳng $\left(AA'M\right) \Rightarrow A'H\perp \left(AB'C'\right)$.\\
		Vậy khoảng cách từ $A'$ đến mặt phẳng $\left(AA'M\right)$ là $A'H$.\\
		Ta có $\dfrac{1}{A'H^2}=\dfrac{1}{A'A^2}+\dfrac{1}{A'M^2}=\dfrac{1}{4a^2}+\dfrac{1}{a^2}\Rightarrow A'H=\dfrac{2a\sqrt{5}}{5}=2$ với $a=\sqrt{5}$.
	}
\end{ex}
\Closesolutionfile{ans}
\indapan{6}{ans/ans-De1-KQ}